\documentclass[11pt,a4paper]{article}

\usepackage[T1]{fontenc}
\usepackage[utf8]{inputenc}
\usepackage[british]{babel}
\usepackage{lmodern}
\usepackage{microtype}
\usepackage[a4paper,margin=1in]{geometry}
\usepackage{amsmath,amssymb,amsthm,mathtools,mathrsfs}
\usepackage{bm}
\usepackage{enumitem}
\usepackage{booktabs,tabularx,array}
\usepackage{tikz}
\usetikzlibrary{arrows.meta,positioning,fit,calc}
\usepackage{xcolor}
\usepackage{hyperref}
\hypersetup{colorlinks=true,linkcolor=blue!60!black,citecolor=blue!60!black,urlcolor=blue!60!black,hypertexnames=false}
\setlength{\emergencystretch}{2em}

\newtheorem{theorem}{Theorem}
\newtheorem{proposition}[theorem]{Proposition}
\newtheorem{lemma}[theorem]{Lemma}
\newtheorem{corollary}[theorem]{Corollary}
\theoremstyle{definition}
\newtheorem{definition}[theorem]{Definition}
\newtheorem{axiom}[theorem]{Axiom}
\newtheorem{remark}[theorem]{Remark}
\newtheorem{example}[theorem]{Example}

\newcommand{\A}{\mathcal A}
\newcommand{\Oset}{\mathcal O}
\newcommand{\M}{\mathcal M}
\newcommand{\R}{\mathbb R}
\newcommand{\D}{\Delta}
\newcommand{\supp}{\operatorname{supp}}
\newcommand{\Id}{\operatorname{Id}}
\newcommand{\KL}{D_{\mathrm{KL}}}
\newcommand{\Sh}{H_{\mathrm{Sh}}}
\newcommand{\I}{I}
\newcommand{\E}{\mathbb E}
\newcommand{\1}{\mathbf 1}
\newcommand{\st}{\mathrm{s.t.}}
\newcommand{\into}{\hookrightarrow}
\newcommand{\eqdist}{\overset{d}{=}}
\newcommand{\co}{\operatorname{co}}

\title{A Preference for Traceable Agency\thanks{Email: \href{mailto:d.condorelli@gmail.com}{d.condorelli@gmail.com}. Substantial portions of this manuscript were developed with assistance from an autonomous, multi-model large-language-model workflow, including generation, critique, and cross-checking of mathematical arguments. The author reviewed, revised, and independently assessed the foundational assumptions and load-bearing steps of the proofs. All claims, errors, and omissions remain the sole responsibility of the author.}}
\author{Daniele Condorelli\\University of Warwick}
\date{\today}

\begin{document}
\maketitle

\begin{abstract}
This paper gives a behavioural foundation for traceable agency: a preference for outcomes that statistically record which
act was realised.  In each finite environment $P:A\to\Delta(O)$, the primitive is a weak order over lotteries on acts, not
an information functional or a preference over channels.  Cross-environment comparisons are encoded as block-supported
comparisons inside one enlarged environment.  Continuity, monotonicity, invariance, and ordinal composition axioms are
necessary and sufficient for $q\succeq_P q'\Longleftrightarrow I_{q,P}(A;O)\ge I_{q',P}(A;O)$.  Thus traceable agency is
entropy reduction, on a common scale across block comparisons.
\end{abstract}

% ===============================================================
\section*{Introduction}
% ===============================================================

Some choices matter not only through the payoff consequences they produce, but through whether those consequences remain
connected to what the agent did.  An elected official may prefer a reform whose realised consequences can be traced to the
chosen intervention rather than the same distribution of consequences generated by an anonymous budget formula.  A donor
may prefer a named project to a general fund even when the distribution of social returns is held fixed.  In such cases
the relevant distinction is not the marginal distribution of final consequences.  It is the joint distribution of acts and
observations: after the outcome is seen, how much can be learned about which act was realised?

This paper axiomatizes that distinction.  I call the motive \emph{traceable agency}.  Traceable agency is not the causal
magnitude of an act's effect, and it is not ordinary expected utility over final payoffs.  An act may have large material
effects that are absorbed into an aggregate outcome and hence leave little trace; another act may have small material
effects but be perfectly identifiable from what is observed.  The question isolated here is whether observable outcomes
make agency statistically legible.  Attribution, responsibility, expressive value, and reputational concerns are possible
reasons why such legibility may matter, but they are not primitives in the theorem.  The object represented below is the
trace itself: the information in outcomes about acts.

The formal environment is deliberately simple.  There is a finite set $A$ of acts, a finite set $O$ of observable
outcomes, and a stochastic kernel $P:A\to\Delta(O)$.  A procedure $q\in\Delta(A)$ first draws a realised act $A$;
conditional on that act, the environment generates $O$ according to $P$.  The outcome may include payoff consequences,
public records, audit signals, or any other observable consequence.  The primitive behavioural datum, however, is only a
weak order $\succeq_P$ on $\Delta(A)$ for each fixed environment $P$.  Thus $q\succeq_P q'$ means that, in environment
$P$, procedure $q$ displays at least as much traceable agency as procedure $q'$.

The environment is not itself an object of choice.  This is important.  The theory does not begin with a preference over
pairs $(q,P)$, nor with a numerical index assigned to channels.  When two environments must be compared, they are embedded
as labelled blocks of a single larger environment.  A comparison between $q$ under $P$ and $p$ under $Q$ is therefore an
ordinary comparison between two block-supported lotteries in $P\sqcup Q$.  This device keeps the behavioural primitive
within the familiar decision-theoretic form---a ranking of alternatives in a fixed problem---while still allowing the
representation to speak across environments.

The axioms are chosen to express the traceability interpretation directly.  The first group is structural.  Rankings are
weak orders and vary continuously.  Duplicating, adding, or deleting irrelevant blocks does not change a comparison.  The
only strictness requirement in this structural group is local non-triviality: whenever there is genuine uncertainty about
which act is realised, full revelation of the act is strictly more traceable than no revelation.  This is not a hidden
normalization of entropy.  It says only that, for every full-support lottery on a non-singleton action set, an outcome
that identifies the realised act has strictly more traceability than an outcome that never changes beliefs about it.  It
fixes the orientation of the criterion and rules out the degenerate one.

The next two axioms are monotonicity requirements.  If the observed outcome is garbled after it is generated, the act
cannot become more traceable.  If the act variable itself is replaced by a noisier or coarser description, the resulting
procedure cannot display more traceability.  The latter axiom is stated for all completions of the induced coarsened
channel: rows of the coarsened action variable that are reached with zero probability must not affect the comparison, no
matter how they are completed.

The remaining axioms discipline composition.  Public-coin independence says that replacing both sides of a fixed-prior
comparison by the same public mixture with a common background experiment leaves the comparison invariant.  Branchwise
continuation monotonicity says that, after a first-stage observation, branchwise improvements in continuation traceability
should improve the compound procedure, with strictness whenever a positive-probability branch strictly improves.
Independent-background separability says that the ranking of a foreground act--outcome relation should not depend on
which unrelated source of traceability is placed beside it.  These are behavioural invariance and dominance requirements,
not numerical information equations.

The main result is an if-and-only-if characterization.  A family $\{\succeq_P\}_P$ satisfies the axioms if and only if,
for every finite channel $P:A\to\Delta(O)$ and every $q,q'\in\Delta(A)$,
\[
q\succeq_P q'
\quad\Longleftrightarrow\quad
I_{q,P}(A;O)\ge I_{q',P}(A;O).
\]
The cross-environment part of the representation is essential: the same mutual-information scale represents
block-supported comparisons across environments.  Hence the theorem does not merely say that each fixed environment has a
ranking ordinally equivalent to mutual information.  It identifies a common unit of traceability for comparing procedures
carried by different stochastic environments.

The representation has a simple interpretation.  Observing $O$ induces a posterior belief about the realised act.  The
traceability of $q$ in $P$ is the expected reduction in uncertainty about $A$:
\[
I_{q,P}(A;O)
=
H(q)-\E_{q,P}\bigl[H(q(\cdot\mid O))\bigr].
\]
If outcomes are independent of acts, the value is zero.  If the act is known in advance because $q$ is degenerate, the
value is also zero.  If each act in the support of $q$ produces a distinct observable outcome, the value is $H(q)$.  Thus
the criterion is not a taste for randomization as such.  Randomization matters only because it creates uncertainty about
agency for the outcome to resolve.

The theorem characterizes only the traceability component of preference.  In applications, under additional additive
separability assumptions, it can be combined with an ordinary payoff term, for example
\[
V(q,P)=\E_{q,P}[u(O)]+\lambda I_{q,P}(A;O).
\]
Axiomatizing the payoff term, the cardinal tradeoff parameter $\lambda$, and possible interactions between payoff and
traceability would require additional primitives.  The present result supplies the missing behavioural foundation for the
information term: once a traceable-agency component is isolated, its scale is forced to be mutual information.

The proof has three steps.  First, block coherence, post-processing, and a prior-specific Blackwell argument imply that
only the induced law of posterior beliefs about acts matters.  Second, public-coin independence gives affine structure on
posterior-law domains, while product, branch, and support-face arguments align the initially separate scales.  Third,
full revelation defines an entropy function; the derived recursion is Faddeev's, and the entropy is therefore Shannon
entropy up to the choice of information unit.  Mutual information then follows as entropy reduction.

The paper belongs to the decision-theoretic tradition that enriches the object of preference in order to capture motives
that expected utility over final payoffs alone leaves out.  Kreps \cite{Kreps1979} represents a preference for flexibility
over opportunity sets; Gul and Pesendorfer \cite{GulPesendorfer2001} represent temptation and commitment through
preferences over menus; Bartling, Fehr, and Herz \cite{BartlingFehrHerz2014} study the intrinsic value of decision
rights.  The present paper studies a different object.  Decision rights and feasible acts are held fixed.  What varies is
whether the realised outcome statistically records the act.

The relation to axiomatic information theory is also different in kind.  Mathematical characterizations of entropy and
information typically start from a numerical functional and impose structural equations on distributions, random
variables, or information-loss maps \cite{Faddeev1956,BaezFritzLeinster2011}.  Here the primitive is an ordinal
preference relation.  Blackwell's comparison of experiments ranks signals by their instrumental value in downstream
decision problems \cite{Blackwell1953}; here the ranking is direct, over the trace left by acts in outcomes.
Rational-inattention models make information about states a costly input into action \cite{Sims2003,MatejkaMcKay2015}.
Here information about actions encoded in outcomes is the valued object.  Related empowerment and free-energy approaches
often take mutual information, channel capacity, or related variational quantities as primitive objective terms.  This
paper instead derives the mutual-information objective from ordinal restrictions and represents the whole map
$q\mapsto I_{q,P}(A;O)$, not only its maximum.

% ===============================================================
\section*{Environment}
% ===============================================================

Let $A$ be a finite set of \emph{actions} and $O$ a finite set of \emph{outcomes}.  A channel is a stochastic kernel
$P:A\to\Delta(O)$.  For each finite channel $P$, the primitive behavioural datum is a weak order
\[
\succeq_P\quad\text{on}\quad \Delta(A).
\]
The interpretation is that $q\succeq_P q'$ means: in the fixed environment $P$, lottery $q$ displays at least as much
traceability as $q'$.
The environment $P$ is not an object of choice.  A lottery is the procedure being evaluated; the channel is the fixed
observation technology through which the realised act may be recorded by outcomes.

\paragraph{Comparison environments.}
Let $P:A\to\Delta(O)$ and $Q:B\to\Delta(Y)$ be two finite channels.  Write $A^0$ and $B^1$ for disjoint labelled
copies of $A$ and $B$.  The comparison environment
\[
P\sqcup Q:
A^0\sqcup B^1\to\Delta\bigl((O\times\{0\})\sqcup(Y\times\{1\})\bigr)
\]
is defined by
\[
(P\sqcup Q)((o,0)\mid a^0)=P(o\mid a),
\qquad
(P\sqcup Q)((y,1)\mid b^1)=Q(y\mid b),
\]
with zero probability assigned to outcomes in the other block.  For $q\in\Delta(A)$ and $p\in\Delta(B)$, let
$q^0\in\Delta(A^0\sqcup B^1)$ and $p^1\in\Delta(A^0\sqcup B^1)$ denote the corresponding block-supported lotteries.
Comparisons between $q$ under $P$ and $p$ under $Q$ will be written directly as
\[
q^0\succeq_{P\sqcup Q}p^1.
\]
This is an ordinary comparison between two lotteries inside the single fixed environment $P\sqcup Q$.

\paragraph{Finite block environments.}
More generally, let $\{P_i:A_i\to\Delta(O_i)\}_{i\in K}$ be a finite family of channels, with labelled disjoint copies
$A_i^i$ of the action sets.  The block environment
\[
\bigsqcup_{i\in K}P_i:
\bigsqcup_{i\in K}A_i^i
\to
\Delta\left(\bigsqcup_{i\in K}(O_i\times\{i\})\right)
\]
uses channel $P_i$ on block $A_i^i$:
\[
\left(\bigsqcup_{k\in K}P_k\right)((o,i)\mid a_i^i)=P_i(o\mid a_i),
\]
with zero probability assigned to outcomes in other blocks.  For $q_i\in\Delta(A_i)$, write $q_i^i$ for the lottery
supported on block $A_i^i$.
Finite block environments are only a notational extension of the two-block construction: they put several otherwise
separate act--outcome environments into one fixed comparison problem.

\paragraph{Posteriors and posterior law.}
Given $(q,P)$, define the joint law
\[
J_{q,P}(a,o):=q(a)P(o\mid a),
\]
the outcome marginal
\[
m_{q,P}(o):=\sum_{a\in A}q(a)P(o\mid a),
\]
and, whenever $m_{q,P}(o)>0$, the posterior
\[
r_o^{\,q,P}(a):=\frac{q(a)P(o\mid a)}{m_{q,P}(o)}.
\]
The posterior law is
\[
\mu_{q,P}:=
\sum_{o:m_{q,P}(o)>0}m_{q,P}(o)\,\delta_{r_o^{\,q,P}}.
\]
It is the distribution of beliefs about the realised act induced by the observed outcome; it records both which posterior
beliefs can arise and with what probabilities.
It is Bayes-plausible:
\[
\int_{\Delta(A)}r\,d\mu_{q,P}(r)=q.
\]
Let $\mathsf X_A:=\Delta_f(\Delta(A))$ be the finite-support posterior-law space on $A$, and write
\[
b(\mu):=\int_{\Delta(A)}r\,d\mu(r)
\]
for the barycentre of $\mu\in\mathsf X_A$.  For $q\in\Delta(A)$, set
\[
\mathcal M_q:=\{\mu\in\mathsf X_A:b(\mu)=q\}.
\]
For a nonempty subset $B\subseteq A$, write
\[
\Delta_B(A):=\{p\in\Delta(A):\operatorname{supp}(p)\subseteq B\},
\]
and let $\rho_B:\Delta_B(A)\to\Delta(B)$ be coordinate restriction, $\rho_B(p)=p|_B$.  If
$\mu\in\mathsf X_A$ satisfies $\operatorname{supp}(b(\mu))\subseteq B$, then
$\operatorname{supp}(\mu)\subseteq\Delta_B(A)$: for each $a\notin B$,
\[
0=b(\mu)(a)=\int_{\Delta(A)}p(a)\,d\mu(p),
\]
and non-negativity implies $p(a)=0$ for every atom $p$ of $\mu$.  In this case define the support projection
\[
\pi_B(\mu):=(\rho_B)_\#\mu
=\sum_{p\in\operatorname{supp}(\mu)}\mu(\{p\})\,\delta_{p|_B}
\in\mathsf X_B.
\]
It preserves barycentres on the face:
\[
b(\pi_B(\mu))=b(\mu)|_B.
\]
In particular, if $b(\mu)=r$ and $\operatorname{supp}(r)=B$, then $\pi_B(\mu)\in\mathcal M_{r|_B}$.
When the barycentre assigns zero probability outside $B$, all posterior beliefs lie on the face $\Delta_B(A)$; $\pi_B$
identifies those beliefs with distributions on $B$.

\paragraph{Public-coin mixing.}
For channels $P:A\to\Delta(O_P)$ and $Q:A\to\Delta(O_Q)$ and $\lambda\in[0,1]$, define
\[
\lambda P\oplus(1-\lambda)Q:
A\to\Delta\bigl((O_P\times\{0\})\sqcup(O_Q\times\{1\})\bigr)
\]
by
\[
(\lambda P\oplus(1-\lambda)Q)((o,0)\mid a)=\lambda P(o\mid a),
\]
and
\[
(\lambda P\oplus(1-\lambda)Q)((o,1)\mid a)=(1-\lambda)Q(o\mid a).
\]
Then
\[
\mu_{q,\lambda P\oplus(1-\lambda)Q}
=
\lambda\mu_{q,P}+(1-\lambda)\mu_{q,Q}.
\]
Because the coin realisation is observed, the observer knows which experiment was used, and the induced posterior law is
the corresponding convex mixture of posterior laws.

\paragraph{Outcome post-processing.}
For a stochastic kernel $T:O\to\Delta(O')$, define
\[
(PT)(o'\mid a):=\sum_{o\in O}P(o\mid a)T(o'\mid o).
\]
The kernel $T$ is a garbling applied after the original outcome is generated.  It changes the recorded observation while
leaving the realised act and the pre-garbling channel $P$ fixed.

\paragraph{Action pushforward.}
For a stochastic kernel $S:A\to\Delta(A')$, define
\[
(qS)(a'):=\sum_{a\in A}q(a)S(a'\mid a),
\qquad
(rS)(a'):=\sum_{a\in A}r(a)S(a'\mid a).
\]
For $\mu\in\mathsf X_A$, write
\[
\mu S:=\mu\circ(r\mapsto rS)^{-1}\in\mathsf X_{A'}.
\]
The kernel $S$ reports the realised act through a possibly noisy description.  Given $P:A\to\Delta(O)$ and
$q\in\Delta(A)$, the Bayesian pushforward channel is the prior-dependent channel from reported acts to outcomes induced by
the joint law $q(a)S(a'\mid a)P(o\mid a)$.  Formally, define
\[
S^qP:A'\to\Delta(O)
\]
by
\[
S^qP(o\mid a')
:=
\frac{\sum_{a\in A}q(a)S(a'\mid a)P(o\mid a)}{(qS)(a')}
\]
whenever $(qS)(a')>0$.  On rows $a'$ with $(qS)(a')=0$, Bayes' rule imposes no value.  A completion of
$S^qP$ is any channel $\widehat P:A'\to\Delta(O)$ that agrees with the displayed formula on every row with
$(qS)(a')>0$.  For every such completion,
\[
\mu_{qS,\widehat P}=\mu_{q,P}S,
\]
since zero-mass rows receive no probability under $qS$.  When a statement uses $S^qP$ without naming a completion, it is
interpreted for every completion of the Bayesian pushforward.

\paragraph{Two-stage composition.}
Let $P_1:A\to\Delta(O_1)$ be a first-stage channel.  For each $o\in O_1$, let
$Q^o:A\to\Delta(O_2^o)$ be a continuation channel.  Define
\[
P_1\triangleright\{Q^o\}_{o\in O_1}:
A\to\Delta\left(\bigsqcup_{o\in O_1}(\{o\}\times O_2^o)\right)
\]
by
\[
(P_1\triangleright\{Q^o\})(o,z\mid a)=P_1(o\mid a)Q^o(z\mid a).
\]
The first signal selects a posterior branch; the continuation channel $Q^o$ is the observation technology used within that
branch.
If $q$ is the initial lottery, write $m(o):=m_{q,P_1}(o)$ and, for each $o$ with $m(o)>0$,
$r_o:=r_o^{\,q,P_1}$.  Zero-probability branches ($m(o)=0$) do not affect the induced posterior law and are ignored
throughout.

\paragraph{Product environments.}
For $P_i:A_i\to\Delta(O_i)$, define
\[
(P_1\otimes P_2)(o_1,o_2\mid a_1,a_2)
:=
P_1(o_1\mid a_1)P_2(o_2\mid a_2).
\]
Product environments combine statistically independent act--outcome components.
For $q_i\in\Delta(A_i)$, write $q_1\otimes q_2$ for the product lottery.
For posterior laws
\[
\mu=\sum_i m_i\delta_{r_i}\in\mathsf X_{A_1},
\qquad
\rho=\sum_j n_j\delta_{s_j}\in\mathsf X_{A_2},
\]
write
\[
\mu\otimes\rho
:=
\sum_{i,j}m_i n_j\,\delta_{r_i\otimes s_j}
\in\mathsf X_{A_1\times A_2}.
\]

\paragraph{Full revelation and no revelation.}
For $q\in\Delta(A)$, the no-information posterior law is $\delta_q$.  The full-revelation posterior law is
\[
\chi_q:=\sum_{a\in A}q(a)\delta_{e_a},
\]
where $e_a$ is the degenerate lottery on action $a$.
No information leaves the observer's belief at the prior $q$; full revelation identifies the realised action.

% ===============================================================
\section*{Axioms}
% ===============================================================

The axioms are imposed directly on the family $\{\succeq_P\}_P$.  They do not assume a posterior-law order, and they do
not introduce a preference over pairs $(q,P)$.  Cross-environment-looking statements are always written as comparisons
between block-supported lotteries inside one fixed block environment.

The axioms treat traceable agency as a comparative property of an act lottery together with an observation procedure.  They
require stable rankings within each environment, coherent comparisons through labelled block environments, monotonicity
when outcomes or act descriptions become less directly traceable, and disciplined behaviour under public randomisation,
sequential observation, and independent background components.

\begin{description}[leftmargin=*,style=sameline]

\item[(A1) Weak order and local non-triviality.]
For every finite channel $P:A\to\Delta(O)$, $\succeq_P$ is complete and transitive on $\Delta(A)$.  The family is
locally non-trivial: for every finite action set $A$ with $|A|\ge 2$ and every full-support $q\in\Delta(A)$,
\[
q^0\succ_{\Id_A\sqcup U_A}q^1,
\]
where $\Id_A:A\to\Delta(A)$ is full revelation and $U_A:A\to\Delta(\{\ast\})$ is the one-outcome uninformative channel.

This axiom gives each observation environment an ordinary weak order over act lotteries and fixes the orientation of the
comparison.  When the realised act is genuinely uncertain, observing it directly is strictly more traceable than observing
nothing about it.

\item[(A2) Continuity.]
For every pair of finite alphabets $(A,O)$, the set
\[
\{(P,q,q')\in \Delta(O)^A\times\Delta(A)\times\Delta(A):q\succeq_P q'\}
\]
is closed.  Thus within-environment rankings vary continuously with both the lotteries and the fixed channel.
In addition, block comparisons are continuous in posterior-law convergence at a fixed full-support prior: if
$q\in\Delta(A)$ has full support, $P_n,Q_n,P,Q$ are finite channels on $A$,
\[
\mu_{q,P_n}\Rightarrow \mu_{q,P},
\qquad
\mu_{q,Q_n}\Rightarrow \mu_{q,Q},
\]
and
\[
q^0\succeq_{P_n\sqcup Q_n}q^1
\qquad\text{for all }n,
\]
then
\[
q^0\succeq_{P\sqcup Q}q^1.
\]

Continuity says that traceability comparisons are stable under small changes in the procedure, the environment, and the
induced posterior laws.  Small modelling perturbations should not create abrupt reversals.

\item[(A3) Block-comparison coherence.]
First, duplicating the same environment into two labelled blocks does not change the original ranking of lotteries:
for every channel $P:A\to\Delta(O)$ and every $q,q'\in\Delta(A)$,
\[
q\succeq_P q'
\quad\Longleftrightarrow\quad
q^0\succeq_{P\sqcup P}(q')^1.
\]
Second, block comparisons are invariant to adding or deleting irrelevant blocks.  For every finite block environment
$\bigsqcup_{k\in K}P_k$, every two distinct blocks $i,j\in K$, and every
$q_i\in\Delta(A_i)$ and $q_j\in\Delta(A_j)$,
\[
q_i^i\succeq_{\bigsqcup_{k\in K}P_k}q_j^j
\quad\Longleftrightarrow\quad
q_i^0\succeq_{P_i\sqcup P_j}q_j^1.
\]

Labelled blocks are a way to place procedures carried by different channels in a common comparison problem.  The labels
should not themselves create traceability, alter the ranking inside a duplicated environment, or let unused blocks affect
the comparison between active blocks.

\item[(A4) Outcome post-processing aversion.]
Let $P:A\to\Delta(O)$ and let $T:O\to\Delta(O')$ be stochastic.  Then, for every $q\in\Delta(A)$,
\[
q^0\succeq_{P\sqcup PT}q^1.
\]

Post-processing changes only the observed signal.  Since a garbling can only discard or obscure information contained in
the original outcome, the original procedure is weakly more traceable than its post-processed version.

\item[(A5) Action-coarsening aversion.]
Let $S:A\to\Delta(A')$ be stochastic.  For every $P:A\to\Delta(O)$, every $q\in\Delta(A)$, and every completion
$\widehat P:A'\to\Delta(O)$ of the Bayesian pushforward $S^qP$,
\[
q^0\succeq_{P\sqcup \widehat P}(qS)^1.
\]
When $qS$ has full support, the completion is unique and we write it simply as $S^qP$; otherwise, any use of $S^qP$ in this
axiom is understood for every completion.  In particular, bijective relabelings are neutral by applying the axiom in both
directions, using the corresponding completions on zero-probability rows if needed.

The act whose trace is evaluated may be reported through a possibly noisy description.  The induced comparison is then
between traceability of the original act and traceability of the described act generated from it.  Such a description
cannot display more observable trace than the act from which it is generated.  Zero-probability rows of the described act
are irrelevant to the comparison, no matter how they are completed.

\item[(A6) Public-coin independence.]
Fix an action set $A$, a lottery $q\in\Delta(A)$, and channels $P,Q,R$ on $A$.  For every $\lambda\in(0,1)$,
\[
q^0\succeq_{P\sqcup Q}q^1
\quad\Longleftrightarrow\quad
q^0\succeq_{(\lambda P\oplus(1-\lambda)R)\sqcup(\lambda Q\oplus(1-\lambda)R)}q^1.
\]

The public coin selects which experiment is run, and the coin realisation is observed.  At a fixed prior this forms a
convex mixture of posterior laws.  Replacing each side of a comparison by the same public mixture with a common background
experiment therefore leaves the original comparison unchanged.

\item[(A7) Branchwise continuation monotonicity.]
Fix $q\in\Delta(A)$ and a first-stage channel $P_1:A\to\Delta(O_1)$.  Let
$m(o):=m_{q,P_1}(o)$ and $r_o:=r_o^{\,q,P_1}$ for all branches with $m(o)>0$.

For each such branch, let $Q^o,R^o:A\to\Delta(O_2^o)$ be continuation channels.  If
\[
r_o^0\succeq_{Q^o\sqcup R^o}r_o^1
\qquad\text{for every }o\text{ with }m(o)>0,
\]
then
\[
q^0
\succeq_{(P_1\triangleright\{Q^o\})\sqcup(P_1\triangleright\{R^o\})}
q^1.
\]
If one branch comparison is strict and has positive probability, then the aggregate comparison is strict.

Sequential observation is evaluated branch by branch.  After a first-stage signal, if one continuation leaves at least as
much trace as another at every posterior that can be reached, then replacing the continuation branch by branch should not
reduce overall traceability.  If the improvement is strict on a reached branch, the aggregate comparison is strict.

\item[(A8) Independent-background separability.]
Let $P_1,Q_1:A_1\to\Delta(O_1)$ be channels on the first component and let $R_2,S_2$ be any channels on $A_2$.  For all
full-support $q_1\in\Delta(A_1)$ and $q_2\in\Delta(A_2)$,
\[
(q_1\otimes q_2)^0
\succeq_{(P_1\otimes R_2)\sqcup(Q_1\otimes R_2)}
(q_1\otimes q_2)^1
\quad\Longleftrightarrow\quad
(q_1\otimes q_2)^0
\succeq_{(P_1\otimes S_2)\sqcup(Q_1\otimes S_2)}
(q_1\otimes q_2)^1.
\]
The symmetric condition holds for comparisons in the second component while the first component channels are held fixed.

A common product component is statistically separate from the foreground component and is held fixed across the two
alternatives being compared.  It may carry traceability about its own action, but changing this common background should
not reverse the ranking of the foreground act--outcome relation.  The full-support restriction keeps all product actions
active in the comparison.

\end{description}

% ===============================================================
\section*{Benchmark Verification}
% ===============================================================

Before proving the representation theorem, we verify that the benchmark family of preferences generated by mutual
information satisfies the axioms.  This establishes consistency of the axiom system and gives the easy direction of the
characterisation.  The substantive content of the paper is the converse direction, proved in the sufficiency section.

\begin{lemma}[Benchmark: mutual information satisfies the axioms]\label{lem:MIbenchmark}
For every finite channel $P:A\to\Delta(O)$, define
\[
q\succeq^I_P q'
\quad\Longleftrightarrow\quad
I_{q,P}(A;O)\ge I_{q',P}(A;O).
\]
Then the family $\{\succeq^I_P\}_P$ satisfies Axioms~\textup{(A1)--(A8)}.
\end{lemma}

\begin{proof}
Write $\Sh$ for Shannon entropy on finite distributions, with $0\log0:=0$, and abbreviate $I_{q,P}(A;O)$ to
$I_{q,P}$ when the alphabets are clear.  We repeatedly use the elementary identities
\[
I_{q,P}
=
\Sh(m_{q,P})-\sum_{a\in A}q(a)\Sh(P(\cdot\mid a)),
\tag{N1}
\]
and
\[
I_{q,P}
=
\Sh(q)-\int_{\Delta(A)}\Sh(r)\,d\mu_{q,P}(r).
\tag{N2}
\]
Both formulae are read on the support of the relevant distributions; zero-probability rows and outcomes do not affect the
value.

\emph{Axiom \textup{(A1)}.}
For each fixed $P$, the relation $\succeq^I_P$ is complete and transitive because it is represented by the real-valued
function $q\mapsto I_{q,P}$.  If $|A|\ge2$ and $q$ has full support, then full revelation gives
\[
I_{q,\Id_A}=\Sh(q)>0,
\]
whereas the one-outcome uninformative channel $U_A$ gives $I_{q,U_A}=0$.  Hence
\[
q^0\succ^I_{\Id_A\sqcup U_A}q^1.
\]

\emph{Axiom \textup{(A2)}.}
For fixed finite alphabets, (N1) shows that $(P,q)\mapsto I_{q,P}$ is continuous, since Shannon entropy is continuous on
every finite simplex.  Therefore the set
\[
\{(P,q,q'):q\succeq^I_P q'\}
=
\{(P,q,q'):I_{q,P}\ge I_{q',P}\}
\]
is closed.

For the posterior-law continuity clause, let $q$ have full support and suppose
\[
\mu_{q,P_n}\Rightarrow\mu_{q,P},
\qquad
\mu_{q,Q_n}\Rightarrow\mu_{q,Q},
\]
and
\[
q^0\succeq^I_{P_n\sqcup Q_n}q^1
\quad\text{for all }n.
\]
For block-supported lotteries,
\[
I_{q^0,P_n\sqcup Q_n}=I_{q,P_n},
\qquad
I_{q^1,P_n\sqcup Q_n}=I_{q,Q_n}.
\]
Thus $I_{q,P_n}\ge I_{q,Q_n}$.  By (N2), weak convergence of posterior laws, and continuity of $\Sh$ on the compact
simplex $\Delta(A)$,
\[
I_{q,P_n}
=
\Sh(q)-\int\Sh(r)\,d\mu_{q,P_n}(r)
\longrightarrow
\Sh(q)-\int\Sh(r)\,d\mu_{q,P}(r)
=
I_{q,P},
\]
and similarly $I_{q,Q_n}\to I_{q,Q}$.  Hence $I_{q,P}\ge I_{q,Q}$, so
$q^0\succeq^I_{P\sqcup Q}q^1$.

\emph{Axiom \textup{(A3)}.}
For every channel $P$ and lotteries $q,q'\in\Delta(A)$,
\[
I_{q^0,P\sqcup P}=I_{q,P},
\qquad
I_{(q')^1,P\sqcup P}=I_{q',P},
\]
because a block-supported lottery puts probability one on its own block and the block label is constant under that
lottery.  Hence
\[
q\succeq^I_P q'
\quad\Longleftrightarrow\quad
q^0\succeq^I_{P\sqcup P}(q')^1.
\]
The same observation proves invariance to adding or deleting irrelevant blocks.  If $q_i^i$ and $q_j^j$ are supported on
blocks $i$ and $j$, their mutual-information values in $\bigsqcup_{k\in K}P_k$ are exactly
\[
I_{q_i,P_i}
\quad\text{and}\quad
I_{q_j,P_j},
\]
independently of all other blocks.

\emph{Axiom \textup{(A4)}.}
Let $T:O\to\Delta(O')$ be a stochastic outcome post-processing.  Under $PT$, the variables satisfy the Markov chain
\[
A\longrightarrow O\longrightarrow O'.
\]
The data-processing inequality gives $I_{q,P}(A;O)\ge I_{q,PT}(A;O')$.  Therefore
\[
q^0\succeq^I_{P\sqcup PT}q^1.
\]

\emph{Axiom \textup{(A5)}.}
Let $S:A\to\Delta(A')$ be stochastic.  Generate $A\sim q$, then $A'$ from $S(\cdot\mid A)$, and $O$ from
$P(\cdot\mid A)$.  Then $A'$ is a garbling of $A$, and
\[
A'\longleftarrow A\longrightarrow O
\]
is a Markov fork.  Data processing gives $I(A;O)\ge I(A';O)$.  The joint law of $(A',O)$ is exactly the one induced by the
Bayesian pushforward $S^qP$ on every row $a'$ with $(qS)(a')>0$.  Values of a completion $\widehat P$ on rows with
$(qS)(a')=0$ do not affect mutual information.  Therefore
\[
I_{q,P}(A;O)\ge I_{qS,\widehat P}(A';O)
\]
for every completion $\widehat P$ of $S^qP$, and so
\[
q^0\succeq^I_{P\sqcup\widehat P}(qS)^1.
\]

\emph{Axiom \textup{(A6)}.}
For public randomisation, the randomisation label is observed and independent of the action.  Thus
\[
I_{q,\lambda P\oplus(1-\lambda)R}
=
\lambda I_{q,P}+(1-\lambda)I_{q,R},
\]
and likewise
\[
I_{q,\lambda Q\oplus(1-\lambda)R}
=
\lambda I_{q,Q}+(1-\lambda)I_{q,R}.
\]
Since $\lambda>0$,
\[
I_{q,P}\ge I_{q,Q}
\quad\Longleftrightarrow\quad
I_{q,\lambda P\oplus(1-\lambda)R}
\ge
I_{q,\lambda Q\oplus(1-\lambda)R}.
\]
This is exactly Axiom~\textup{(A6)}.

\emph{Axiom \textup{(A7)}.}
Fix a first-stage channel $P_1:A\to\Delta(O_1)$.  Let
\[
m(o)=m_{q,P_1}(o),
\qquad
r_o=r_o^{\,q,P_1}.
\]
For a continuation profile $\mathcal Q=\{Q^o\}_o$, the chain rule for mutual information gives
\[
I_{q,P_1\triangleright\mathcal Q}(A;O_1,O_2)
=
I_{q,P_1}(A;O_1)
+
\sum_{o:m(o)>0}m(o)\,I_{r_o,Q^o}(A;O_2^o).
\tag{N3}
\]
Therefore, if
\[
r_o^0\succeq^I_{Q^o\sqcup R^o}r_o^1
\quad\text{for every positive-probability branch }o,
\]
then $I_{r_o,Q^o}\ge I_{r_o,R^o}$ for every such $o$.  Using (N3),
\[
I_{q,P_1\triangleright\{Q^o\}}
\ge
I_{q,P_1\triangleright\{R^o\}},
\]
so
\[
q^0
\succeq^I_{(P_1\triangleright\{Q^o\})\sqcup(P_1\triangleright\{R^o\})}
q^1.
\]
If one branch comparison is strict and that branch has positive probability, then the corresponding summand in (N3) is
strictly larger, so the aggregate comparison is strict.

\emph{Axiom \textup{(A8)}.}
For product priors and product channels,
\[
I_{q_1\otimes q_2,P_1\otimes R_2}(A_1,A_2;O_1,O_2)
=
I_{q_1,P_1}(A_1;O_1)+I_{q_2,R_2}(A_2;O_2).
\tag{N4}
\]
Consequently,
\[
I_{q_1\otimes q_2,P_1\otimes R_2}
\ge
I_{q_1\otimes q_2,Q_1\otimes R_2}
\]
is equivalent to $I_{q_1,P_1}\ge I_{q_1,Q_1}$.  Replacing $R_2$ by any $S_2$ leaves this last inequality unchanged, proving
the first-component clause of Axiom~\textup{(A8)}.  The symmetric second-component clause is identical.
\end{proof}

% ===============================================================
\section*{Main Result}
% ===============================================================

\begin{theorem}[Mutual-information characterisation]\label{thm:main}
For a family $\{\succeq_P\}_P$ of weak orders over finite stochastic environments, the following are equivalent.
\begin{enumerate}[label=\textup{(\roman*)},leftmargin=*]
\item The family satisfies Axioms~\textup{(A1)--(A8)}.
\item For every finite channel $P:A\to\Delta(O)$ and every $q,q'\in\Delta(A)$,
\[
q\succeq_P q'
\quad\Longleftrightarrow\quad
I_{q,P}(A;O)\ge I_{q',P}(A;O).
\]
\end{enumerate}
Moreover, under these equivalent conditions, block-supported cross-channel comparisons are represented on the same
scale: for every finite block environment $\bigsqcup_{k\in K}P_k$, every two distinct blocks $i,j\in K$, and every
$q_i\in\Delta(A_i)$, $q_j\in\Delta(A_j)$,
\[
q_i^i\succeq_{\bigsqcup_{k\in K}P_k}q_j^j
\quad\Longleftrightarrow\quad
I_{q_i,P_i}(A_i;O_i)\ge I_{q_j,P_j}(A_j;O_j).
\]
\end{theorem}

\begin{proof}
If \textup{(ii)} holds, then $\{\succeq_P\}_P$ is the mutual-information family, so
Lemma~\ref{lem:MIbenchmark} gives \textup{(i)}.  The sufficiency section proves \textup{(i)$\Rightarrow$(ii)},
culminating in Lemma~\ref{lem:globalsketch}.  For the moreover clause, apply
\textup{(ii)} to the fixed block environment $\bigsqcup_{k\in K}P_k$.  A lottery supported on block $i$ has constant
block label, so
\[
I_{q_i^i,\bigsqcup_{k\in K}P_k}=I_{q_i,P_i}(A_i;O_i),
\]
and similarly for block $j$.
\end{proof}

% ===============================================================
\section*{Proof of Sufficiency}
% ===============================================================

Assume Axioms~\textup{(A1)--(A8)}.  We prove the implication \textup{(i)$\Rightarrow$(ii)} in
Theorem~\ref{thm:main}.  The proof is written first for full-support lotteries.  Boundary cases are obtained by
restricting to the support and then using continuity.  The key point is that posterior-law sufficiency is not assumed; it
is derived from block comparisons and outcome post-processing.

\paragraph{Proof roadmap.}
Lemmas~\ref{lem:blockcoh}--\ref{lem:blackwell} derive posterior-law sufficiency from block comparisons and Blackwell
equivalence.  Lemmas~\ref{lem:plsuff}--\ref{lem:actionbase} give fibrewise affine posterior-separable representations
and basic invariances.  Lemma~\ref{lem:neutrality} records undetectable-distinctions neutrality as a derived consequence
of action coarsening.  Lemma~\ref{lem:coherentnorm} chooses a coherent product gauge and proves quasi-additivity, allowing
a possible bilinear interaction term.  Lemma~\ref{lem:blockbridge} moves cross-prior block comparisons into these
unscaled values.  Lemma~\ref{lem:branchagg} derives branch aggregation, and Lemma~\ref{lem:chain} derives the cocycle,
the beta-form, and the normalised chain rule.  Lemma~\ref{lem:facescales} aligns the chain scales across support faces.
Lemma~\ref{lem:scalecoherence} uses the chain rule and a two-grouping argument to eliminate the interaction term and prove
the scale $a_q$ is universal.
Lemma~\ref{lem:faddeevsketch} then proves entropy reduction, Faddeev recursion, and the mutual-information formula.
Lemma~\ref{lem:globalsketch} gives the fixed-environment representation.

\begin{lemma}[Coherent block comparisons]\label{lem:blockcoh}
Under \textup{(A1)} and \textup{(A3)}, pairwise comparisons between block-supported lotteries are independent of the block
environment in which the two blocks are embedded.  In particular, fix a finite action set $A$ and a full-support
$q\in\Delta(A)$.  For channels $P,Q,R$ on $A$, the comparisons
\[
q^0\succeq_{P\sqcup Q}q^1
\]
are complete and transitive as comparisons of channels at prior $q$.  The same Axiom~\textup{(A3)} convention identifies
the reversed block-supported comparison with its canonical two-block form:
\[
q^1\succeq_{P\sqcup Q}q^0
\quad\Longleftrightarrow\quad
q^0\succeq_{Q\sqcup P}q^1.
\]
Under \textup{(A2)}, these comparisons are also closed in the channels: if
$P_n\to P$, $Q_n\to Q$, and
\[
q^0\succeq_{P_n\sqcup Q_n}q^1
\quad\text{for all }n,
\]
then
\[
q^0\succeq_{P\sqcup Q}q^1.
\]
\end{lemma}

\begin{proof}
Completeness follows from Axiom~\textup{(A1)} applied to the two-block environment $P\sqcup Q$: either
$q^0\succeq_{P\sqcup Q}q^1$ or $q^1\succeq_{P\sqcup Q}q^0$.  In the second case, applying
Axiom~\textup{(A3)} to the two-block environment with $P_0=P$ and $P_1=Q$, using the ordered pair $(i,j)=(1,0)$,
gives the canonical reverse comparison
\[
q^1\succeq_{P\sqcup Q}q^0
\quad\Longleftrightarrow\quad
q^0\succeq_{Q\sqcup P}q^1.
\]

For transitivity, suppose
\[
q^0\succeq_{P\sqcup Q}q^1
\qquad\text{and}\qquad
q^0\succeq_{Q\sqcup R}q^1.
\]
Form the three-block environment $P^0\sqcup Q^1\sqcup R^2$, and write $q^P,q^Q,q^R$ for the lottery $q$ supported on
the $P$, $Q$, and $R$ blocks.  By the irrelevant-block clause of Axiom~\textup{(A3)} with block pair $(0,1)$, the first
comparison is equivalent to
\[
q^P\succeq_{P^0\sqcup Q^1\sqcup R^2}q^Q,
\]
and by the same clause with block pair $(1,2)$, the second comparison is equivalent to
\[
q^Q\succeq_{P^0\sqcup Q^1\sqcup R^2}q^R.
\]
Transitivity of the weak order in this common environment gives
\[
q^P\succeq_{P^0\sqcup Q^1\sqcup R^2}q^R.
\]
Applying Axiom~\textup{(A3)} again, now deleting the irrelevant $Q$ block, yields
\[
q^0\succeq_{P\sqcup R}q^1.
\]

For closedness, the channel $P_n\sqcup Q_n$ converges to $P\sqcup Q$ in the finite-dimensional product simplex.  Since
the block-supported lotteries $q^0,q^1$ are fixed, Axiom~\textup{(A2)} gives
\[
q^0\succeq_{P\sqcup Q}q^1.
\]
\end{proof}

\begin{lemma}[Prior-specific Blackwell equivalence]\label{lem:blackwell}
Let $q\in\Delta(A)$ have full support.  Let $P:A\to\Delta(O)$ and $Q:A\to\Delta(Y)$ be finite channels.  If
\[
\mu_{q,P}=\mu_{q,Q},
\]
then there are stochastic kernels $T:O\to\Delta(Y)$ and $T':Y\to\Delta(O)$ such that
\[
Q=PT,
\qquad
P=QT'.
\]
\end{lemma}

\begin{proof}
It is enough to construct $T$; the construction of $T'$ is symmetric.  Let
\[
m_P(o):=m_{q,P}(o),\qquad m_Q(y):=m_{q,Q}(y).
\]
If $m_P(o)=0$, then $P(o\mid a)=0$ for every $a\in A$, because $q$ has full support.  Such outcomes do not affect
$PT$, so define the corresponding row of $T$ arbitrarily.  The same observation applies to null outcomes of $Q$.

Let $\mathcal R$ be the finite set of posterior vectors that occur with positive probability under either channel.  For
$r\in\mathcal R$, define
\[
O_r:=\{o\in O:m_P(o)>0,\ r_o^{\,q,P}=r\},
\qquad
Y_r:=\{y\in Y:m_Q(y)>0,\ r_y^{\,q,Q}=r\}.
\]
The sets $\{O_r\}_{r\in\mathcal R}$ partition the positive-probability outcomes of $P$, and the sets
$\{Y_r\}_{r\in\mathcal R}$ partition the positive-probability outcomes of $Q$, by posterior vector.
Equality of posterior laws implies that the total probability attached to posterior $r$ is the same for the two
channels:
\[
\lambda_r:=\sum_{o\in O_r}m_P(o)=\sum_{y\in Y_r}m_Q(y)>0.
\]
For $o\in O_r$ and $y\in Y_r$, set
\[
T(y\mid o):=\frac{m_Q(y)}{\lambda_r},
\]
and set $T(y\mid o)=0$ when $o\in O_r$ and $y\notin Y_r$.  This defines a stochastic row for every positive-probability
outcome $o$, since the probabilities over $Y_r$ sum to one.  Rows corresponding to null outcomes of $P$ were already
chosen arbitrarily.

Fix $a\in A$ and $y\in Y_r$.  Since all outcomes in $O_r$ have posterior $r$, Bayes' rule gives
\[
q(a)P(o\mid a)=m_P(o)r(a)
\qquad(o\in O_r).
\]
Therefore
\[
\sum_{o\in O_r}P(o\mid a)
=
\frac{\lambda_r r(a)}{q(a)}.
\]
Using the definition of $T$,
\[
(PT)(y\mid a)
=
\sum_{o\in O_r}P(o\mid a)\frac{m_Q(y)}{\lambda_r}
=
\frac{m_Q(y)r(a)}{q(a)}.
\]
But $y\in Y_r$ also has posterior $r$, so Bayes' rule for $Q$ gives
\[
Q(y\mid a)=\frac{m_Q(y)r(a)}{q(a)}.
\]
Thus $(PT)(y\mid a)=Q(y\mid a)$ for every positive-probability $y$.  If $m_Q(y)=0$, then full support of $q$ implies
$Q(y\mid a)=0$ for every $a$.  No positive-probability outcome of $P$ sends mass to such a $y$ under the construction;
rows of $T$ indexed by null outcomes of $P$ may be chosen arbitrarily, but those outcomes have $P(o\mid a)=0$ for every
$a$, so they contribute nothing to $PT$.  Hence $(PT)(y\mid a)=0=Q(y\mid a)$ for every $a$, and therefore $Q=PT$.
\end{proof}

\begin{lemma}[Posterior-law sufficiency]\label{lem:plsuff}
Fix a full-support $q\in\Delta(A)$.  If $P:A\to\Delta(O)$ and $Q:A\to\Delta(Y)$ generate the same posterior law at $q$,
\[
\mu_{q,P}=\mu_{q,Q},
\]
then
\[
q^0\sim_{P\sqcup Q}q^1.
\]
Consequently, by Lemma~\ref{lem:blockcoh}, at a fixed full-support prior the block comparison of channels depends only
on the induced posterior law.
\end{lemma}

\begin{proof}
By Lemma~\ref{lem:blackwell}, there are stochastic kernels $T$ and $T'$ such that $Q=PT$ and $P=QT'$.  Axiom~\textup{(A4)}
gives
\[
q^0\succeq_{P\sqcup PT}q^1,
\qquad
q^0\succeq_{Q\sqcup QT'}q^1.
\]
Using $PT=Q$, the first comparison is
\[
q^0\succeq_{P\sqcup Q}q^1.
\]
Using $QT'=P$, the second comparison is
\[
q^0\succeq_{Q\sqcup P}q^1.
\]
By the reverse-block equivalence in Lemma~\ref{lem:blockcoh} (equivalently, Axiom~\textup{(A3)} applied to
$P^0\sqcup Q^1$ with ordered pair $(i,j)=(1,0)$), this is equivalent to
\[
q^1\succeq_{P\sqcup Q}q^0.
\]
Thus $q^0\sim_{P\sqcup Q}q^1$.
\end{proof}

\begin{lemma}[Posterior-separable representation]\label{lem:postsep}
Fix a full-support $q\in\Delta(A)$.  There is a continuous affine functional
\[
F_q:\mathcal M_q\to\R
\]
such that, for any two channels $P,Q$ on $A$,
\[
q^0\succeq_{P\sqcup Q}q^1
\quad\Longleftrightarrow\quad
F_q(\mu_{q,P})\ge F_q(\mu_{q,Q}).
\]
Moreover, for some continuous $\phi_q:\Delta(A)\to\R$,
\[
F_q(\mu)=\int_{\Delta(A)}\phi_q(r)\,d\mu(r),
\]
and therefore
\[
F_q(\mu_{q,P})
=
\sum_{o:m_{q,P}(o)>0}m_{q,P}(o)\,\phi_q(r_o^{\,q,P}).
\]
The affine functional is unique up to positive affine transformations.  The integrand $\phi_q$ is not pointwise unique;
adding an affine function of $r$ leaves its integral on the fibre $\mathcal M_q$ changed only by a constant.
\end{lemma}

\begin{proof}
First note that every finite-support Bayes-plausible law is implemented by a finite channel.  If
\[
\mu=\sum_{k=1}^n\lambda_k\delta_{r_k}\in\mathcal M_q,
\]
define
\[
P(k\mid a)=\frac{\lambda_k r_k(a)}{q(a)}.
\]
Since $q$ has full support and $\sum_k\lambda_k r_k=q$, this is a channel.  Its outcome probability at $k$ is
$\lambda_k$, and the posterior after $k$ is $r_k$; hence $\mu_{q,P}=\mu$.

Define a comparison on $\mathcal M_q$ as follows: for $\mu,\nu\in\mathcal M_q$, choose any channels $P,Q$ implementing
them and compare the two block-supported copies of $q$ in $P\sqcup Q$.  This comparison is independent of the chosen
implementations.  Indeed, if $P,P'$ both implement $\mu$ and $Q,Q'$ both implement $\nu$, then
Lemma~\ref{lem:plsuff} gives
\[
q^0\sim_{P\sqcup P'}q^1,
\qquad
q^0\sim_{Q\sqcup Q'}q^1.
\]
Placing the relevant blocks in common block environments and using Lemma~\ref{lem:blockcoh}, transitivity gives
\[
q^0\succeq_{P\sqcup Q}q^1
\quad\Longleftrightarrow\quad
q^0\succeq_{P'\sqcup Q'}q^1.
\]
Thus the quotient comparison on $\mathcal M_q$ is well defined; denote it by $\succeq_q$.  This construction is only at
the fixed full-support prior $q$, and every law in $\mathcal M_q$ is finite-support and, by the preceding construction,
implemented by a finite channel.  Lemma~\ref{lem:blockcoh} also gives completeness and transitivity of this quotient
comparison.

The set $\mathcal M_q$ is a mixture space under ordinary convex combination.  It is convex because barycentres are
preserved by convex combination.  The quotient relation $\succeq_q$ is a weak order by Lemma~\ref{lem:blockcoh}.
Public-coin mixtures of implementing channels correspond exactly to convex mixtures of posterior laws:
\[
\mu_{q,\lambda P\oplus(1-\lambda)R}
=
\lambda\mu_{q,P}+(1-\lambda)\mu_{q,R}.
\]
Therefore Axiom~\textup{(A6)} gives mixture independence: for every $\mu,\nu,\rho\in\mathcal M_q$ and
$\lambda\in(0,1)$,
\[
\mu\succeq_q\nu
\quad\Longleftrightarrow\quad
\lambda\mu+(1-\lambda)\rho
\succeq_q
\lambda\nu+(1-\lambda)\rho.
\]

The relative weak topology on $\mathcal M_q$ is inherited from the weak topology on probability measures over the compact
metric simplex $\Delta(A)$, and convex-mixture paths are continuous in this topology.  The fixed-prior posterior-law
continuity clause of Axiom~\textup{(A2)} gives relatively closed upper and lower contour sets.  Indeed, if
$\mu_n\Rightarrow\mu$ with $\mu_n,\mu,z\in\mathcal M_q$ and $\mu_n\succeq_q z$, choose finite implementing channels for
these laws and apply Axiom~\textup{(A2)} to get $\mu\succeq_q z$; the argument for $\{\,\mu:z\succeq_q\mu\,\}$ is the
same.  Therefore, for fixed $\mu,\nu,\rho\in\mathcal M_q$, the sets
\[
\{\lambda\in[0,1]:\lambda\mu+(1-\lambda)\nu\succeq_q\rho\}
\quad\text{and}\quad
\{\lambda\in[0,1]:\rho\succeq_q\lambda\mu+(1-\lambda)\nu\}
\]
are closed.

By the Herstein--Milnor mixture-space theorem \cite{HersteinMilnor1953}, there is a mixture-preserving representation
$F_q:\mathcal M_q\to\R$:
\[
\mu\succeq_q\nu
\quad\Longleftrightarrow\quad
F_q(\mu)\ge F_q(\nu),
\]
and
\[
F_q(\lambda\mu+(1-\lambda)\nu)
=
\lambda F_q(\mu)+(1-\lambda)F_q(\nu).
\]
If the quotient order is nonconstant, this representative is unique up to positive affine transformations.  If the
quotient order is constant, take $F_q\equiv0$.

Moreover, $F_q$ is continuous in the relative weak topology.  For any $z\in\mathcal M_q$, the sets
\[
\{\mu:F_q(\mu)\ge F_q(z)\}
=
\{\mu:\mu\succeq_q z\}
\quad\text{and}\quad
\{\mu:F_q(\mu)\le F_q(z)\}
=
\{\mu:z\succeq_q\mu\}
\]
are relatively closed by the preceding Axiom~\textup{(A2)} argument.  Since $F_q$ is affine, its range on the convex set
$\mathcal M_q$ is an interval.  Hence every upper and lower level set of $F_q$ is relatively closed: levels outside the
range are trivial, and levels inside the range equal $F_q(z)$ for some $z\in\mathcal M_q$.  Thus $F_q$ is both upper and
lower semicontinuous, and therefore continuous.  For any channels $P,Q$ on $A$,
\[
q^0\succeq_{P\sqcup Q}q^1
\quad\Longleftrightarrow\quad
F_q(\mu_{q,P})\ge F_q(\mu_{q,Q}).
\]

It remains to write the affine functional in posterior-separable form, without leaving the finite-support domain.  Let
$X:=\Delta(A)$.  Choose
\[
0<\varepsilon<\min_{a\in A}q(a).
\]
For each $r\in X$, define
\[
s_r:=\frac{q-\varepsilon r}{1-\varepsilon}\in\Delta(A)
\]
and
\[
\eta_r
:=
\varepsilon\delta_r+(1-\varepsilon)\sum_{a\in A}s_r(a)\delta_{e_a}.
\]
Then $\eta_r\in\mathcal M_q$, because its barycentre is
\[
\varepsilon r+(1-\varepsilon)s_r=q.
\]
Choose constants $c_a$ such that
\[
\sum_{a\in A}q(a)c_a=F_q(\chi_q),
\]
for instance $c_a=F_q(\chi_q)$ for every $a$.  Define
\[
\phi_q(r):=
\frac{
F_q(\eta_r)
-
(1-\varepsilon)\sum_{a\in A}s_r(a)c_a
}{\varepsilon}.
\]
The map $r\mapsto\eta_r$ is weakly continuous and $F_q$ is continuous, so $\phi_q$ is continuous.

Now take any finite-support law
\[
\mu=\sum_{k=1}^n\lambda_k\delta_{r_k}\in\mathcal M_q.
\]
Since $\sum_k\lambda_k r_k=q$, we have
\[
\sum_{k=1}^n\lambda_k s_{r_k}=q.
\]
Therefore
\[
\sum_{k=1}^n\lambda_k\eta_{r_k}
=
\varepsilon\mu+(1-\varepsilon)\chi_q.
\]
By affinity of $F_q$,
\[
\sum_{k=1}^n\lambda_k F_q(\eta_{r_k})
=
F_q\!\left(\varepsilon\mu+(1-\varepsilon)\chi_q\right)
=
\varepsilon F_q(\mu)+(1-\varepsilon)F_q(\chi_q).
\]
Using $\sum_k\lambda_k s_{r_k}=q$ and $\sum_a q(a)c_a=F_q(\chi_q)$ gives
\[
\sum_{k=1}^n\lambda_k\phi_q(r_k)=F_q(\mu).
\]
This is exactly
\[
F_q(\mu)=\int_X\phi_q(r)\,d\mu(r)
\]
for every finite-support $\mu\in\mathcal M_q$.  Applying this to $\mu_{q,P}$ gives the stated finite-sum formula.  The
affine representative $F_q$ is unique up to positive affine transformations by Herstein--Milnor.  The integrand is not
pointwise unique on a fixed barycentre fibre: if $\ell$ is affine in $r$, then $\int\ell(r)\,d\mu(r)=\ell(q)$ for every
$\mu\in\mathcal M_q$.
\end{proof}

\begin{lemma}[Refinement monotonicity, convexity, and normalisation]\label{lem:convnorm}
Fix a full-support $q\in\Delta(A)$ and let $F_q,\phi_q$ be as in Lemma~\ref{lem:postsep}.  If a finite law
$\mu\in\mathcal M_q$ is obtained from another finite law $\nu\in\mathcal M_q$ by replacing each posterior of $\nu$ by a
finite Bayes-plausible continuation experiment, then
\[
F_q(\mu)\ge F_q(\nu).
\]
Consequently, $\phi_q$ may be chosen convex.  After replacing $\phi_q$ by an observationally equivalent representative
and subtracting a constant from $F_q$, one may impose
\[
\phi_q(q)=F_q(\delta_q)=0,
\qquad
\phi_q(r)\ge 0\quad\text{for every }r\in\Delta(A),
\]
and therefore
\[
F_q(\mu)\ge 0\quad\text{for every }\mu\in\mathcal M_q.
\]
\end{lemma}

\begin{proof}
Write
\[
\nu=\sum_{i=1}^n m_i\delta_{r_i},
\qquad
\mu=\sum_{i=1}^n m_i\nu_i,
\qquad
\nu_i\in\mathcal M_{r_i}.
\]
Choose a channel $P:A\to\Delta(\{1,\dots,n\})$ implementing $\nu$ at prior $q$.  For each $i$, choose a finite
continuation channel $Q^i:A\to\Delta(Z_i)$ implementing $\nu_i$ at prior $r_i$; on actions outside the support of
$r_i$, define the rows arbitrarily.  Those rows are never reached under prior $r_i$ and do not affect $\nu_i$; moreover,
if $r_i(a)=0$, then $q(a)P(i\mid a)=m_i r_i(a)=0$, so they do not affect the compound posterior law at $q$.  The compound
channel
\[
C:=P\triangleright\{Q^i\}_{i=1}^n
\]
implements $\mu$ at prior $q$.  Let $T:\bigsqcup_{i=1}^n(\{i\}\times Z_i)\to\Delta(\{1,\dots,n\})$ be the kernel
$T(i\mid(i,z))=1$; then $CT=P$.  Axiom~\textup{(A4)} applied to $C$ and $T$ gives
\[
q^0\succeq_{C\sqcup P}q^1.
\]
By Lemma~\ref{lem:postsep},
\[
F_q(\mu)\ge F_q(\nu).
\]

To prove convexity, take $r,s_1,\dots,s_k\in\Delta(A)$ with
\[
r=\sum_{j=1}^k\lambda_j s_j,
\qquad
\lambda_j\ge 0,\quad \sum_j\lambda_j=1.
\]
Choose $0<\varepsilon<\min_a q(a)$ and set
\[
t:=\frac{q-\varepsilon r}{1-\varepsilon}.
\]
This belongs to $\Delta(A)$: for each $a$, $q(a)-\varepsilon r(a)\ge q(a)-\varepsilon>0$, and the coordinates sum to one.
The law
\[
\varepsilon\sum_{j=1}^k\lambda_j\delta_{s_j}+(1-\varepsilon)\delta_t
\]
is a refinement of
\[
\varepsilon\delta_r+(1-\varepsilon)\delta_t,
\]
and both belong to $\mathcal M_q$.  Refinement monotonicity and the integral representation give
\[
\varepsilon\sum_{j=1}^k\lambda_j\phi_q(s_j)+(1-\varepsilon)\phi_q(t)
\ge
\varepsilon\phi_q(r)+(1-\varepsilon)\phi_q(t).
\]
Cancelling the common term and dividing by $\varepsilon$ yields
\[
\phi_q(r)\le \sum_{j=1}^k\lambda_j\phi_q(s_j),
\]
so $\phi_q$ is convex.

Since $q$ is in the relative interior of the finite-dimensional simplex, the supporting-hyperplane theorem for finite
convex functions gives a relative subgradient $g\in\R^A$ of the continuous convex function $\phi_q$ at $q$.  Define
\[
\widetilde\phi_q(r):=\phi_q(r)-g\cdot(r-q).
\]
For every $\mu\in\mathcal M_q$,
\[
\int g\cdot(r-q)\,d\mu(r)=g\cdot(b(\mu)-q)=0,
\]
so $\widetilde\phi_q$ represents the same affine functional on $\mathcal M_q$.  The subgradient inequality gives
\[
\widetilde\phi_q(r)\ge \widetilde\phi_q(q)
\qquad\text{for all }r\in\Delta(A).
\]
Finally replace
\[
F_q(\mu)\quad\text{by}\quad F_q(\mu)-F_q(\delta_q)
\]
and replace $\widetilde\phi_q$ by
\[
\widetilde\phi_q-\widetilde\phi_q(q).
\]
Before this final constant shift, the representation with $\widetilde\phi_q$ gives
\[
F_q(\delta_q)=\int \widetilde\phi_q(r)\,d\delta_q(r)=\widetilde\phi_q(q).
\]
Thus the same constant is subtracted from both sides of the integral representation, and subtracting a constant from
$F_q$ preserves all ordinal comparisons on $\mathcal M_q$.  The final normalisation gives
\[
\phi_q(q)=F_q(\delta_q)=0,
\]
and makes $\phi_q$ non-negative on the whole simplex.  Hence
\[
F_q(\mu)=\int\phi_q(r)\,d\mu(r)\ge 0
\]
for every $\mu\in\mathcal M_q$.
\end{proof}

\begin{lemma}[Cross-fibre ordinal equivalence]\label{lem:actionbase}
Action coarsening implies ordinal cross-fibre equivalence for bijections.  In particular, all no-information
experiments have the same normalised value zero, and for every bijection $\pi:A\to A'$ and full-support
$q\in\Delta(A)$, there exists $c_{\pi,q}>0$ such that
\[
F_{q\pi}(\mu\pi)=c_{\pi,q} F_q(\mu)
\]
for all $\mu\in\mathcal M_q$.  The exact equality $c_{\pi,q}=1$ for every $q$ will be derived as
Corollary~\ref{cor:permutationinvariance} after product normalisation is established.
\end{lemma}

\begin{proof}
\emph{Ordinal equivalence.}  Let $\pi:A\to A'$ be a bijection.  View $\pi$ as the deterministic kernel
$S(a'\mid a)=\1_{a'=\pi(a)}$.  For any channel $P:A\to\Delta(O)$ and full-support $q\in\Delta(A)$, the Bayesian
pushforward satisfies $S^qP=P\pi^{-1}$ (explicitly: $(S^qP)(o\mid a')=P(o\mid\pi^{-1}(a'))$).  Axiom~\textup{(A5)}
gives
\[
q^0\succeq_{P\sqcup P\pi^{-1}}(q\pi)^1.
\]
For the reverse inequality, apply Axiom~\textup{(A5)} to the bijection $\pi^{-1}:A'\to A$, the prior $q\pi\in\Delta(A')$,
and the channel $P\pi^{-1}:A'\to\Delta(O)$.  Then $(\pi^{-1})^{q\pi}(P\pi^{-1})=P$ (the inverse bijection undoes the
relabeling), and $(q\pi)\pi^{-1}=q$.  Hence
\[
(q\pi)^0\succeq_{P\pi^{-1}\sqcup P}q^1.
\]
By the second part of Axiom~\textup{(A3)}, block comparisons depend only on the pair of channels involved, not on
numeric labels.  The comparison $(q\pi)^0\succeq_{P\pi^{-1}\sqcup P}q^1$ asks whether $(q\pi,P\pi^{-1})$ is at least
as good as $(q,P)$.  Swapping the block labels (so that $P$ becomes block~0 and $P\pi^{-1}$ becomes block~1) gives
the equivalent statement $q^0\preceq_{P\sqcup P\pi^{-1}}(q\pi)^1$.  Combining with the first inequality gives
\[
q^0\sim_{P\sqcup P\pi^{-1}}(q\pi)^1.
\]
If $P$ implements $\mu\in\mathcal M_q$, then $P\pi^{-1}$ implements $\mu\pi\in\mathcal M_{q\pi}$: for every
positive-probability outcome $o$,
\[
m_{q\pi,P\pi^{-1}}(o)=m_{q,P}(o),
\qquad
r_o^{\,q\pi,P\pi^{-1}}=r_o^{\,q,P}\pi.
\]
Now take arbitrary $\mu,\nu\in\mathcal M_q$, and choose channels $P,Q$ implementing them at prior $q$.  Let $E$ be the
four-block environment with blocks $P$, $P\pi^{-1}$, $Q$, and $Q\pi^{-1}$.  Write $x,x',y,y'$ for the block-supported
lotteries $q$, $q\pi$, $q$, and $q\pi$ on these four blocks, respectively.  By Lemma~\ref{lem:blockcoh}, the two
two-block indifferences give $x\sim_E x'$ and $y\sim_E y'$.  Hence, by transitivity of the weak order on $E$,
\[
x\succeq_E y
\quad\Longleftrightarrow\quad
x'\succeq_E y'.
\]
Applying Lemma~\ref{lem:blockcoh} again to delete irrelevant blocks gives
\[
q^0\succeq_{P\sqcup Q}q^1
\quad\Longleftrightarrow\quad
(q\pi)^0\succeq_{P\pi^{-1}\sqcup Q\pi^{-1}}(q\pi)^1.
\]
By Lemma~\ref{lem:postsep} on the two fibres, this is equivalent to
\[
F_q(\mu)\ge F_q(\nu)
\quad\Longleftrightarrow\quad
F_{q\pi}(\mu\pi)\ge F_{q\pi}(\nu\pi).
\]
The map $\mu\mapsto\mu\pi$ is an affine bijection from $\mathcal M_q$ to $\mathcal M_{q\pi}$, with inverse induced by
$\pi^{-1}$.  Therefore $G(\mu):=F_{q\pi}(\mu\pi)$ is a continuous affine representative of the same weak order on
$\mathcal M_q$ as $F_q$.  Herstein--Milnor uniqueness gives $G=aF_q+b$ with $a>0$.  Since
$\delta_q\pi=\delta_{q\pi}$ and Lemma~\ref{lem:convnorm} gives
\[
F_q(\delta_q)=F_{q\pi}(\delta_{q\pi})=0,
\]
we have $b=0$.  Setting $c_{\pi,q}:=a$ yields
\[
F_{q\pi}(\mu\pi)=c_{\pi,q} F_q(\mu)
\]
for all $\mu\in\mathcal M_q$.

\emph{No-information equivalence.}  By Lemma~\ref{lem:convnorm}, $F_q(\delta_q)=0$ for every full-support $q$ on every
finite action set, independently of cardinality.
\end{proof}

\begin{lemma}[Undetectable distinctions neutrality]\label{lem:neutrality}
Let $S:A\to A'$ be an onto deterministic partition map, and let $P:A\to\Delta(O)$ be $S$-measurable:
\[
S(a)=S(b)\quad\Rightarrow\quad P(\cdot\mid a)=P(\cdot\mid b).
\]
Then for every full-support $q\in\Delta(A)$,
\[
q^0\sim_{P\sqcup S^qP}(qS)^1.
\]
In particular, if $G_S:A\to\Delta(A')$ reveals the block $S(a)$, then
\[
(q,G_S)\sim(qS,\Id_{A'}).
\]
\end{lemma}

\begin{proof}
Since $S$ is onto and $q$ has full support, $qS$ has full support on $A'$.  Thus the Bayesian pushforward $S^qP$ has no
zero-row ambiguity in this lemma.  View the deterministic map $S$ as its associated stochastic kernel.
One direction is exactly Axiom~\textup{(A5)}:
\[
q^0\succeq_{P\sqcup S^qP}(qS)^1.
\]
For the reverse direction, define the Bayesian refinement kernel $R:A'\to\Delta(A)$ by
\[
R(a\mid a')=
\begin{cases}
q(a)/(qS)(a'),&S(a)=a',\\
0,&\text{otherwise}.
\end{cases}
\]
For each $a'\in A'$, $R(\cdot\mid a')$ is stochastic because
\[
\sum_{a\in A}R(a\mid a')
=
\frac{\sum_{a:S(a)=a'}q(a)}{(qS)(a')}
=1.
\]
Moreover, for every $a\in A$,
\[
((qS)R)(a)
=
\sum_{a'\in A'}(qS)(a')R(a\mid a')
=
(qS)(S(a))\frac{q(a)}{(qS)(S(a))}
=q(a)>0.
\]
Thus $(qS)R=q$.

Since $P$ is $S$-measurable, the row of $S^qP$ at each block $a'$ equals the common $P$-row on that block: if $S(a)=a'$,
then
\[
S^qP(o\mid a')
=
\frac{\sum_{b:S(b)=a'}q(b)P(o\mid b)}{(qS)(a')}
=
P(o\mid a).
\]
Now pushing $S^qP$ back through $R$ recovers $P$.  For each $a\in A$,
\[
\bigl(R^{qS}(S^qP)\bigr)(o\mid a)
=
\frac{\sum_{a'\in A'}(qS)(a')R(a\mid a')S^qP(o\mid a')}{((qS)R)(a)}
=
S^qP(o\mid S(a))
=
P(o\mid a).
\]
Therefore
\[
R^{qS}(S^qP)=P.
\]

Apply Axiom~\textup{(A5)} to $(qS,S^qP)$ with the refinement kernel $R$:
\[
(qS)^0\succeq_{S^qP\sqcup P}q^1.
\]
By Axiom~\textup{(A3)} applied to the two-block environment $P\sqcup S^qP$ with ordered pair $(1,0)$, this is equivalent
to
\[
(qS)^1\succeq_{P\sqcup S^qP}q^0.
\]
Together the two inequalities give the desired indifference.
\end{proof}

\begin{lemma}[Coherent product quasi-additivity]\label{lem:coherentnorm}
There is a choice of positive rescalings of the zero-normalised fibrewise representatives $F_q$, still denoted $F_q$,
and a scalar $\kappa\in\R$ such that, for all full-support priors $p\in\Delta(A)$ and $r\in\Delta(B)$ and all
$\mu\in\mathcal M_p$, $\nu\in\mathcal M_r$,
\begin{equation}
F_{p\otimes r}(\mu\otimes\nu)
=F_p(\mu)+F_r(\nu)+\kappa F_p(\mu)F_r(\nu).
\tag{QA}\label{eq:QA}
\end{equation}
If one factor is a singleton, the same formula holds with its identically zero representative.  Moreover, for every
nondegenerate $r$ and every $\nu\in\mathcal M_r$,
\begin{equation}
1+\kappa F_r(\nu)>0.
\tag{POS}\label{eq:QAslope}
\end{equation}
\end{lemma}

\begin{proof}
All representatives in this proof are zero-normalised: $F_q(\delta_q)=0$.  The proof first treats nondegenerate action
sets, meaning action sets with at least two elements.

\emph{Step 1: product-coordinate independence.}
Fix full-support priors $p$ and $r$ on nondegenerate finite action sets and set
\[
L_{p,r}(\mu,\nu):=F_{p\otimes r}(\mu\otimes\nu).
\]
This map is continuous and separately affine.  To identify a first-coordinate slice, let $P,Q$ implement
$\mu,\mu'\in\mathcal M_p$ and let $R$ implement $\nu\in\mathcal M_r$.  Axiom~\textup{(A8)} gives
\[
(p\otimes r)^0\succeq_{(P\otimes R)\sqcup(Q\otimes R)}(p\otimes r)^1
\quad\Longleftrightarrow\quad
(p\otimes r)^0\succeq_{(P\otimes U_B)\sqcup(Q\otimes U_B)}(p\otimes r)^1.
\]
Let $S:A\times B\to A$ be the projection.  Then $(p\otimes r)S=p$, and the Bayesian pushforward of
$P\otimes U_B$ under $S$ is $P$ up to the trivial one-outcome relabelling of the background channel; the same holds for
$Q\otimes U_B$ and $Q$.  Projection neutrality (Lemma~\ref{lem:neutrality}) therefore identifies
$(p\otimes r,P\otimes U_B)$ with $(p,P)$ and $(p\otimes r,Q\otimes U_B)$ with $(p,Q)$ in block comparisons.  These two
indifferences and the comparison above are placed in a common block environment using Lemma~\ref{lem:blockcoh};
transitivity then gives
\[
L_{p,r}(\mu,\nu)\ge L_{p,r}(\mu',\nu)
\quad\Longleftrightarrow\quad
F_p(\mu)\ge F_p(\mu').
\]
The symmetric argument identifies every second-coordinate slice with the $F_r$-order.  Local non-triviality, the
fibrewise representation, and the zero normalisation from Lemma~\ref{lem:convnorm} make both slice orders nonconstant:
$F_p(\chi_p)>0$ and $F_r(\chi_r)>0$.

\emph{Step 2: the pair-specific bilinear form.}
Put $x(\mu):=F_p(\mu)$ and $y(\nu):=F_r(\nu)$.  For fixed $\nu$, affine-utility uniqueness gives
\[
L_{p,r}(\mu,\nu)=\alpha(\nu)x(\mu)+\gamma(\nu),
\qquad \alpha(\nu)>0.
\]
Taking $\mu=\delta_p$ shows $\gamma(\nu)=L_{p,r}(\delta_p,\nu)$.  This is a continuous affine representative of the
$F_r$-order and vanishes at $\delta_r$, so
\[
\gamma(\nu)=B_{p,r}y(\nu)
\]
for some $B_{p,r}>0$.  Likewise, the slice at $\nu=\delta_r$ gives
\[
L_{p,r}(\mu,\delta_r)=A_{p,r}x(\mu)
\]
for some $A_{p,r}>0$.

Let $\bar\mu:=\chi_p$ and $h_p:=F_p(\chi_p)>0$.  The function
$\nu\mapsto L_{p,r}(\bar\mu,\nu)$ is again a continuous affine representative of the $F_r$-order.  Hence there are
$D_{p,r}>0$ and a constant $E_{p,r}$ such that
\[
L_{p,r}(\bar\mu,\nu)=D_{p,r}y(\nu)+E_{p,r}.
\]
At $\nu=\delta_r$, $E_{p,r}=A_{p,r}h_p$.  Comparing this identity with
$L_{p,r}(\bar\mu,\nu)=\alpha(\nu)h_p+B_{p,r}y(\nu)$ gives
\[
\alpha(\nu)=A_{p,r}+C_{p,r}y(\nu),
\qquad
C_{p,r}:=\frac{D_{p,r}-B_{p,r}}{h_p}.
\]
Therefore
\begin{equation}
L_{p,r}(\mu,\nu)
=A_{p,r}x(\mu)+B_{p,r}y(\nu)+C_{p,r}x(\mu)y(\nu).
\tag{BIL}\label{eq:pairbilinear}
\end{equation}
The coordinate-slice slopes are positive:
\[
A_{p,r}+C_{p,r}y(\nu)>0,
\qquad
B_{p,r}+C_{p,r}x(\mu)>0.
\]

\emph{Step 3: normalising the two linear coefficients.}
Adopt the strict-product convention: finite products are identified by tuple notation, so
$(A\times B)\times C=A\times B\times C=A\times(B\times C)$ and similarly for product priors.  This convention does
not identify $A\times B$ with $B\times A$.

Expanding $F_{p\otimes r\otimes s}$ under the two groupings and comparing the coefficients of $x$, $y$, and $z$ gives
\begin{align}
A_{p\otimes r,s}A_{p,r}&=A_{p,r\otimes s}, \tag{C1}\label{eq:Acoc1}\\
A_{p\otimes r,s}B_{p,r}&=B_{p,r\otimes s}A_{r,s}, \tag{C2}\label{eq:Acoc2}\\
B_{p\otimes r,s}&=B_{p,r\otimes s}B_{r,s}. \tag{C3}\label{eq:Acoc3}
\end{align}
Coefficient comparison is legitimate because the mixtures $(1-t)\delta_p+t\chi_p$ make $x$ range over a nondegenerate
interval, and similarly in the other coordinates.

The coordinate swap gives a scalar $d(p,r)>0$ such that
\[
F_{r\otimes p}(\nu\otimes\mu)=d(p,r)F_{p\otimes r}(\mu\otimes\nu).
\]
Comparing coefficients in~\eqref{eq:pairbilinear},
\[
B_{r,p}=d(p,r)A_{p,r},
\qquad
A_{r,p}=d(p,r)B_{p,r},
\qquad
C_{r,p}=d(p,r)C_{p,r}.
\]
Thus, with $\rho(p,r):=B_{p,r}/A_{p,r}$,
\begin{equation}
\rho(r,p)=\rho(p,r)^{-1}.
\tag{SW}\label{eq:rhoswap}
\end{equation}

Divide~\eqref{eq:Acoc2} by~\eqref{eq:Acoc1}:
\[
\rho(p,r)=\rho(p,r\otimes s)A_{r,s}.
\]
Fix a nondegenerate reference prior $q_0$ on a finite action set $A_0$ and put $\phi(r):=\rho(q_0,r)$.  Then
\begin{equation}
A_{r,s}=\frac{\phi(r)}{\phi(r\otimes s)}.
\tag{G}\label{eq:gaugeA}
\end{equation}
Rescale every representative by $F_p^*:=\phi(p)F_p$.  Under this positive rescaling,
\[
A^*_{p,r}=A_{p,r}\frac{\phi(p\otimes r)}{\phi(p)},
\quad
B^*_{p,r}=B_{p,r}\frac{\phi(p\otimes r)}{\phi(r)},
\quad
C^*_{p,r}=C_{p,r}\frac{\phi(p\otimes r)}{\phi(p)\phi(r)}.
\]
Equation~\eqref{eq:gaugeA} gives $A^*_{p,r}=1$ for every pair.  Write $\beta(p,r):=B^*_{p,r}$.  The transformed
versions of~\eqref{eq:Acoc2}--\eqref{eq:Acoc3} and~\eqref{eq:rhoswap} are
\begin{align*}
\beta(p,r)&=\beta(p,r\otimes s),\\
\beta(p\otimes r,s)&=\beta(p,r)\beta(r,s),\\
\beta(r,p)&=\beta(p,r)^{-1}.
\end{align*}
Moreover $\phi(q_0)=\rho(q_0,q_0)=1$.  Equation~\eqref{eq:gaugeA} gives
$A_{q_0,p}=1/\phi(q_0\otimes p)$, while
$B_{q_0,p}=A_{q_0,p}\rho(q_0,p)=\phi(p)/\phi(q_0\otimes p)$.  Therefore
\[
\beta(q_0,p)=B_{q_0,p}\frac{\phi(q_0\otimes p)}{\phi(p)}=1.
\]
By swap reciprocity, $\beta(p,q_0)=1$.  Taking $s=q_0$ in the second displayed identity gives
\[
1=\beta(p\otimes r,q_0)=\beta(p,r)\beta(r,q_0)=\beta(p,r).
\]
Thus, after the positive gauge rescaling, $A_{p,r}=B_{p,r}=1$ for every nondegenerate pair.  Positivity of the slice
slopes is preserved because the rescaling factors are positive.

\emph{Step 4: the interaction coefficient is universal.}
Write the remaining coefficient as $\kappa_{p,r}$.  Comparing the $xy$, $xz$, $yz$, and $xyz$ coefficients in the two
associative expansions gives
\begin{align}
\kappa_{p,r}&=\kappa_{p,r\otimes s}, \tag{K1}\label{eq:kappa1}\\
\kappa_{p\otimes r,s}&=\kappa_{p,r\otimes s}, \tag{K2}\label{eq:kappa2}\\
\kappa_{p\otimes r,s}&=\kappa_{r,s}, \tag{K3}\label{eq:kappa3}\\
\kappa_{p\otimes r,s}\kappa_{p,r}&=\kappa_{p,r\otimes s}\kappa_{r,s}. \tag{K4}\label{eq:kappa4}
\end{align}
The last equation follows from the first three, but it records the full coefficient comparison.  After $A=B=1$, the
linear coefficient comparison under the coordinate swap forces the swap scalar to be one; the bilinear coefficient
comparison then gives
\begin{equation}
\kappa_{p,r}=\kappa_{r,p}.
\tag{KS}\label{eq:kappasym}
\end{equation}
Using~\eqref{eq:kappa1}--\eqref{eq:kappa3} with the reference prior $q_0$,
\[
\kappa_{p,r}
=\kappa_{p,r\otimes q_0}
=\kappa_{p\otimes r,q_0}
=\kappa_{r,q_0}.
\]
Hence $\kappa_{p,r}$ is independent of its first argument.  By symmetry it is also independent of its second argument.
Denote the common value by $\kappa$.  This proves~\eqref{eq:QA} for nondegenerate factors.

The positive slice-slope condition now reads
\[
1+\kappa F_r(\nu)>0,
\]
which is~\eqref{eq:QAslope}.

\emph{Step 5: singleton factors.}
Let $q_*=\delta_*$ be the prior on a singleton.  Its zero-normalised representative is identically zero.  If $p$ is
nondegenerate, the canonical bijection $A\to A\times\{*\}$ and Lemma~\ref{lem:actionbase} give a scalar $c>0$ with
\[
F_{p\otimes q_*}(\mu\otimes\delta_{q_*})=cF_p(\mu).
\]
Tensor both sides with the nondegenerate reference prior $q_0$.  The canonical bijection between
$(A\times\{*\})\times A_0$ and $A\times A_0$ gives a scalar $d>0$, and quasi-additivity for the two nondegenerate
products gives, with $x=F_p(\mu)$ and $y=F_{q_0}(\nu)$,
\[
cx+y+\kappa cxy=d(x+y+\kappa xy).
\]
Taking $\mu=\delta_p$ and $\nu=\chi_{q_0}$ gives $x=0$ and $y>0$, hence $d=1$.  Taking
$\mu=\chi_p$ and $\nu=\delta_{q_0}$ gives $x>0$ and $y=0$, hence $c=1$.  Thus~\eqref{eq:QA} also holds with a singleton
factor.  The reversed order is identical.
\end{proof}

\begin{corollary}[Exact relabelling invariance]\label{cor:permutationinvariance}
For every bijection $\pi:A\to A'$ and every full-support $q\in\Delta(A)$,
\[
F_{q\pi}(\mu\pi)=F_q(\mu)
\qquad(\mu\in\mathcal M_q).
\]
\end{corollary}

\begin{proof}
The singleton case is immediate because both representatives vanish.  Suppose $A$ is nondegenerate.  Lemma~\ref{lem:actionbase}
gives $F_{q\pi}(\mu\pi)=cF_q(\mu)$ for some $c>0$.  Apply the product bijection
$\pi\times\mathrm{id}_{A_0}$ with the nondegenerate reference prior $q_0$.  There is $d>0$ such that, writing
$x=F_q(\mu)$ and $y=F_{q_0}(\nu)$,
\[
cx+y+\kappa cxy=d(x+y+\kappa xy).
\]
Set $\mu=\delta_q$ and $\nu=\chi_{q_0}$, so $x=0$ and $y>0$, to obtain $d=1$.  Then set
$\mu=\chi_q$ and $\nu=\delta_{q_0}$, so $x>0$ and $y=0$, to obtain $c=1$.
\end{proof}

\begin{lemma}[Support restriction]\label{lem:supprestrict}
Let $r\in\Delta(A)$ have support $B:=\operatorname{supp}(r)\subsetneq A$.  Then:
\begin{enumerate}
\item[(i)] \emph{Posterior-law reduction.}  For any channel $P:A\to\Delta(O)$, the posterior law $\mu_{r,P}$ depends only
on the rows $P(\cdot\mid b)$ for $b\in B$.  Rows $P(\cdot\mid a)$ with $a\notin B$ have zero contribution because $r(a)=0$.
\item[(ii)] \emph{Preference equivalence.}  Two channels $P,Q:A\to\Delta(O)$ that agree on $B$ (i.e.\ $P(\cdot\mid b)=Q(\cdot\mid b)$
for all $b\in B$) are indifferent at prior $r$.
\item[(iii)] \emph{Face identification.}  The comparison between channels $P,Q$ at prior $r$ agrees with the comparison
between $P|_B,Q|_B$ at the full-support prior $r|_B\in\Delta(B)$.  This is proved by projecting $A$ to $B$ and embedding
$B$ back into $A$, using the all-completions form of Axiom~\textup{(A5)} for the embedding.
\item[(iv)] \emph{Representative.}  Define $F_r:\mathcal M_r\to\R$ by $F_r(\mu):=F_{r|_B}^B(\pi_B(\mu))$, where $\pi_B$
projects posterior laws from $\Delta(A)$ to $\Delta(B)$ and $F_{r|_B}^B$ is the Lemma~\ref{lem:postsep} representative
on $\Delta(B)$.  This $F_r$ represents continuation comparisons at prior $r$ in action set $A$.
\end{enumerate}
\end{lemma}

\begin{proof}
In this proof, $(s,R)\succeq(t,T)$ abbreviates the block comparison
$s^0\succeq_{R\sqcup T}t^1$, and $\sim$ denotes both weak inequalities.  Whenever several such comparisons are chained,
we first place the relevant blocks in a common finite block environment using Axiom~\textup{(A3)} and then apply
Axiom~\textup{(A1)} transitivity.

Part~(i): The marginal probability $m_{r,P}(o)=\sum_a r(a)P(o\mid a)$ and the Bayes numerator $r(a)P(o\mid a)$ are both
zero for $a\notin B$.  Hence $\mu_{r,P}$ is determined by $\{P(\cdot\mid b):b\in B\}$.

Part~(ii): Let $P$ and $\widetilde P$ be channels on $A$ that agree on $B$.  We first show
$(r,P)\sim(r,\widetilde P)$.  Fix $b_0\in B$, and let $S:A\to B$ be the deterministic map with $S(b)=b$ for $b\in B$
and $S(a)=b_0$ for $a\notin B$.  Then $rS=r|_B$ and the Bayesian pushforward $S^rP$ is $P|_B$ on $B$.  Axiom~\textup{(A5)}
gives
\[
(r,P)\succeq(r|_B,P|_B).
\]
Let $I:B\hookrightarrow A$ be the inclusion.  Then $(r|_B)I=r$, and the Bayesian pushforward $I^{r|_B}(P|_B)$ is
pinned down on $B$ and unconstrained on $A\setminus B$.  By the all-completions convention in Axiom~\textup{(A5)}, choose
the completion to be $\widetilde P$.  Axiom~\textup{(A5)} gives
\[
(r|_B,P|_B)\succeq(r,\widetilde P).
\]
Thus $(r,P)\succeq(r,\widetilde P)$.  Applying the same projection/inclusion argument with $\widetilde P$ as the original
channel and choosing $P$ as the inclusion completion gives $(r,\widetilde P)\succeq(r,P)$.  Hence
$(r,P)\sim(r,\widetilde P)$.

If $P,Q$ agree on $B$, choose a common extension $\widetilde P=\widetilde Q$ of this common restriction to all of $A$
(for instance, use one fixed off-support row).  The preceding paragraph gives
$(r,P)\sim(r,\widetilde P)$ and $(r,Q)\sim(r,\widetilde Q)$; since $\widetilde P=\widetilde Q$, common-environment
transitivity gives $(r,P)\sim(r,Q)$.

Part~(iii): We first prove $(r,P)\sim(r|_B,P|_B)$ for every channel $P:A\to\Delta(O)$.

\emph{Forward direction (A5 projection).}  Let $S:A\to B$ map all $a\notin B$ to some fixed $b_0\in B$.  The coarsened
prior is $rS=r|_B$ (since $r(a)=0$ for $a\notin B$), and the Bayesian pushforward $S^rP$ equals $P|_B$ on $B$.  By
Axiom~\textup{(A5)},
\[
(r,P)\succeq(r|_B,P|_B).
\]

\emph{Reverse direction (A5 inclusion with completion).}  Let $I:B\hookrightarrow A$ be the inclusion.  Then
$(r|_B)I=r$.  The Bayesian pushforward $I^{r|_B}(P|_B)$ is pinned down on $B$ and arbitrary on $A\setminus B$.  Choose
the Axiom~\textup{(A5)} completion to be any extension $\widetilde P:A\to\Delta(O)$ of $P|_B$.  Then
\[
(r|_B,P|_B)\succeq(r,\widetilde P).
\]
By Part~(ii), $(r,\widetilde P)\sim(r,P)$ because $\widetilde P$ and $P$ agree on $B=\operatorname{supp}(r)$.  Hence
$(r|_B,P|_B)\succeq(r,P)$.

Together the two directions give $(r,P)\sim(r|_B,P|_B)$.

\emph{Comparison equivalence.}  For any $P,Q$: $(r,P)\succeq(r,Q)$ iff (by the above indifferences $(r,P)\sim(r|_B,P|_B)$,
$(r,Q)\sim(r|_B,Q|_B)$, and Axiom~\textup{(A1)} transitivity) $(r|_B,P|_B)\succeq(r|_B,Q|_B)$.

Part~(iv): If $\mu\in\mathcal M_r$, every posterior in the finite support of $\mu$ is supported on $B$; otherwise the
barycentre $r$ would assign positive mass to $A\setminus B$.  Hence $\pi_B(\mu)\in\mathcal M_{r|_B}$ is well defined.
By (iii), continuation comparisons at $r$ are isomorphic to those at $r|_B$.  The Lemma~\ref{lem:postsep} representative
$F_{r|_B}^B$ on $\Delta(B)$ therefore induces a representative on $\mathcal M_r$ via $\pi_B$.
\end{proof}

\begin{lemma}[Cross-prior block representation]\label{lem:blockbridge}
For arbitrary finite priors $q\in\Delta(A)$ and $r\in\Delta(B)$ and channels
$P:A\to\Delta(O)$ and $Q:B\to\Delta(Y)$,
\begin{equation}
q^0\succeq_{P\sqcup Q}r^1
\quad\Longleftrightarrow\quad
F_q(\mu_{q,P})\ge F_r(\mu_{r,Q}).
\tag{BC}\label{eq:blockcomp}
\end{equation}
Boundary representatives are interpreted by Lemma~\ref{lem:supprestrict}.
\end{lemma}

\begin{proof}
First suppose $q$ and $r$ have full support, allowing singleton action sets.  At prior $q\otimes r$, quasi-additivity and
zero normalisation give
\[
F_{q\otimes r}(\mu_{q,P}\otimes\delta_r)=F_q(\mu_{q,P}),
\qquad
F_{q\otimes r}(\delta_q\otimes\mu_{r,Q})=F_r(\mu_{r,Q}).
\]
These values are in the same fibre $\mathcal M_{q\otimes r}$, so Lemma~\ref{lem:postsep} gives
\[
(q\otimes r)^0\succeq_{(P\otimes U_B)\sqcup(U_A\otimes Q)}(q\otimes r)^1
\quad\Longleftrightarrow\quad
F_q(\mu_{q,P})\ge F_r(\mu_{r,Q}).
\]

It remains to transfer the product-lifted comparison back to the original two blocks.  For the $A$-side pair,
Axiom~\textup{(A5)} applied to the projection $\pi_A:A\times B\to A$ gives
\[
(q\otimes r,P\otimes U_B)\succeq(q,P)
\]
in block comparisons.  Conversely, let $T:A\to\Delta(A\times B)$ be
\[
T((a,b)\mid a')=\1_{\{a=a'\}}r(b).
\]
Then $qT=q\otimes r$, and the Bayesian pushforward $T^qP$ has row $P(\cdot\mid a)$ at every $(a,b)$.  Thus
$T^qP$ is $P\otimes U_B$ after the canonical one-outcome identification of the $B$-experiment, and Axiom~\textup{(A5)}
gives the reverse comparison.  Equivalently, if the harmless outcome identification is not made, Lemma~\ref{lem:plsuff}
identifies $T^qP$ and $P\otimes U_B$ because they induce the same posterior law at $q\otimes r$.  Hence
$(q\otimes r,P\otimes U_B)$ and $(q,P)$ are indifferent in every relevant block comparison.  The same
projection/embedding argument identifies $(q\otimes r,U_A\otimes Q)$ with $(r,Q)$.

Place the four pairs $(q,P)$, $(q\otimes r,P\otimes U_B)$, $(r,Q)$, and $(q\otimes r,U_A\otimes Q)$ as blocks of one
finite block environment.  Lemma~\ref{lem:blockcoh} and transitivity transfer the two pairwise indifferences to the
two-block comparison, proving~\eqref{eq:blockcomp} for full-support priors.

For arbitrary priors, let $A_+:=\operatorname{supp}(q)$ and $B_+:=\operatorname{supp}(r)$.  Lemma~\ref{lem:supprestrict}
identifies $(q,P)$ with $(q|_{A_+},P|_{A_+})$, identifies $(r,Q)$ with $(r|_{B_+},Q|_{B_+})$, and identifies the two
values $F_q(\mu_{q,P})$ and $F_r(\mu_{r,Q})$ with their support-face values.  Applying the full-support bridge on
$A_+$ and $B_+$ and pulling back by Lemma~\ref{lem:supprestrict} proves the claim.  Singleton supports use the same
zero-representative convention as Lemma~\ref{lem:coherentnorm}.
\end{proof}

\begin{lemma}[Branch aggregation]\label{lem:branchagg}
Let $P_1:A\to\Delta(O_1)$ be a first-stage channel at full-support $q$, with branch probabilities $m(o)$ and posteriors
$r_o$.  There are positive branch coefficients $\beta(q,r_o)$, depending only on $q$ and the reached posterior $r_o$,
such that for every continuation profile $\{Q^o\}_{o:m(o)>0}$,
\[
F_q(\mu_{q,P_1\triangleright\{Q^o\}})
=
F_q(\mu_{q,P_1})+
\sum_{o:m(o)>0}m(o)\,\beta(q,r_o)\,
F_{r_o}(\mu_{r_o,Q^o}).
\]
If $r_o$ has full support, $\beta(q,r_o)$ is the full-support tangent coefficient.  If
$1<|\operatorname{supp}(r_o)|<|A|$, it is the full-to-face coefficient.  If $|\operatorname{supp}(r_o)|=1$, the value
$F_{r_o}$ is identically zero and $\beta(q,r_o)$ is an arbitrary positive convention.
\end{lemma}

\begin{proof}
Fix the first-stage channel $P_1$ and its branch structure $(m(o),r_o)_{o:m(o)>0}$.  For any profile of continuation
channels $\{Q^o\}$, the compound posterior law at prior $q$ is
\[
\mu_{q,P_1\triangleright\{Q^o\}}
=
\sum_{o:m(o)>0}m(o)\,\mu_{r_o,Q^o}.
\]

\emph{Setup: linear part of affine functional.}  Since $F_q$ is affine on $\mathcal M_q$ (Lemma~\ref{lem:postsep}),
it extends uniquely to an affine functional on the affine hull of $\mathcal M_q$.  Write $L_q$ for the linear part:
for any $\mu,\mu'\in\mathcal M_q$,
\[
F_q(\mu)-F_q(\mu')=L_q(\mu-\mu').
\]

\emph{Branch-affine structure.}  Fix a positive-probability branch $\bar o$ and fix continuation channels
$H^o:A\to\Delta(Z^o)$ in all other branches $o\ne\bar o$.  Let $r:=r_{\bar o}$ and $m:=m(\bar o)$.  If
$\operatorname{supp}(r)$ is a singleton, then $\mathcal M_r=\{\delta_r\}$ and $F_r(\delta_r)=0$ by the singleton
zero-normalisation convention; such branches contribute zero and their coefficient convention is recorded in
Case~1 below.  For the branch-affine argument in this paragraph, assume $|\operatorname{supp}(r)|\ge2$.  If $r$ is a
boundary posterior, $F_r$ and continuation comparisons at $r$ are understood through Lemma~\ref{lem:supprestrict}.

Let
\[
\rho_H:=\sum_{o\ne\bar o}m(o)\,\mu_{r_o,H^o}.
\]
For any $\nu\in\mathcal M_r$, the aggregate posterior law produced by using branch law $\nu$ in branch $\bar o$ and
the fixed continuations elsewhere is
\[
T_H(\nu):=m\nu+\rho_H\in\mathcal M_q.
\]
Define $g:\mathcal M_r\to\R$ by
\[
g(\nu):=F_q(T_H(\nu)).
\]
Thus $g$ depends only on $\nu$, and $g=F_q\circ T_H$ is continuous and affine.

We now show that $g$ and $F_r$ represent the same weak order on $\mathcal M_r$.  For $\nu,\nu'\in\mathcal M_r$, choose
finite continuation channels implementing them at prior $r$; when $r$ is boundary these are implemented on
$B=\operatorname{supp}(r)$ and extended arbitrarily to $A$, as licensed by Lemma~\ref{lem:supprestrict}.  Padding outcome
alphabets if needed, Axiom~\textup{(A7)} compares the two continuation profiles that differ only in branch $\bar o$.
If $F_r(\nu)>F_r(\nu')$, Axiom~\textup{(A7)} gives $g(\nu)>g(\nu')$.  Conversely, if $g(\nu)\ge g(\nu')$ but
$F_r(\nu)<F_r(\nu')$, applying the strict part of Axiom~\textup{(A7)} to the swapped branch comparison gives
$g(\nu')>g(\nu)$, a contradiction.  Hence $g(\nu)\ge g(\nu')$ iff $F_r(\nu)\ge F_r(\nu')$.

Since $|\operatorname{supp}(r)|\ge2$, local non-triviality on the support face and Lemmas~\ref{lem:supprestrict} and
\ref{lem:convnorm} imply $F_r$ is nonconstant.  The continuous affine functions $g$ and $F_r$ are therefore nonconstant
affine representatives of the same weak order on $\mathcal M_r$.  Herstein--Milnor uniqueness gives
\[
g(\nu)=\alpha\,F_r(\nu)+\gamma
\]
for some $\alpha>0$ and $\gamma\in\R$.

\emph{The coefficient $\alpha$ is independent of the fixed continuations.}  Write $g^{(1)}$ and $g^{(2)}$ for the
functionals arising from two different choices $\{H^o\}$ and $\{\widetilde H^o\}$ in the other branches.  Both
represent the same order as $F_r$, so $g^{(1)}=\alpha_1 F_r+\gamma_1$ and
$g^{(2)}=\alpha_2 F_r+\gamma_2$ with $\alpha_1,\alpha_2>0$.  To show $\alpha_1=\alpha_2$, consider the
difference $g^{(1)}-g^{(2)}=(\alpha_1-\alpha_2)F_r+(\gamma_1-\gamma_2)$.  For any
$\nu\in\mathcal M_r$,
\[
g^{(1)}(\nu)-g^{(2)}(\nu)
=
F_q\!\left(m(\bar o)\nu+\textstyle\sum_{o\ne\bar o}m(o)\mu_{r_o,H^o}\right)
-
F_q\!\left(m(\bar o)\nu+\textstyle\sum_{o\ne\bar o}m(o)\mu_{r_o,\widetilde H^o}\right).
\]
Both arguments are posterior laws in $\mathcal M_q$, differing only in the fixed non-$\bar o$ components.  Since $L_q$
has already been defined on the affine-hull difference space,
\[
g^{(1)}(\nu)-g^{(2)}(\nu)=L_q(\rho_1-\rho_2),
\]
where $\rho_i$ is the fixed non-$\bar o$ background law for $g^{(i)}$.  This difference is independent of $\nu$.  Since
$F_r$ is nonconstant in the present nondegenerate case, $\alpha_1=\alpha_2$.

\emph{Path-independence of the branch coefficient via tangent-space decomposition.}
We now prove that $\alpha/m(\bar o)$ depends only on the pair $(q,r_{\bar o})$, not on the first-stage experiment.
This is the key step that converts ordinal branchwise monotonicity into a cardinal aggregation formula.

\emph{Tangent space at a posterior.}  For a posterior $r\in\Delta(A)$ with support $B:=\operatorname{supp}(r)$, define
the \emph{tangent space}:
\[
V_r:=\operatorname{span}\{\nu-\nu':\nu,\nu'\in\mathcal M_r\}.
\]
Elements of $V_r$ are signed measures with total mass zero and barycentre zero.

\emph{Tangent-space identification.}  For any $r$ with full support, we claim
\[
V_r=W:=\{\eta\in\operatorname{span}\Delta_f(\Delta(A)):\eta(\Delta(A))=0,\ \textstyle\int p\,d\eta(p)=0\},
\]
the space of finite-support signed measures on $\Delta(A)$ with zero total mass and zero barycentre.
Indeed, if $\nu,\nu'\in\mathcal M_r$, then $\nu-\nu'$ has total mass $1-1=0$ and barycentre $r-r=0$, so $V_r\subset W$.
Conversely, given any $\eta\in W$: if $\eta=0$ then $\eta\in V_r$ trivially.  Otherwise, since $\eta$ has finite support,
	write its Jordan decomposition $\eta=\eta^+-\eta^-$ where $\eta^+,\eta^-\ge 0$ are finite-support positive measures
	with equal total mass $c>0$.  Normalise $\zeta^+:=\eta^+/c$ and $\zeta^-:=\eta^-/c$; both are finite-support probability
	measures on $\Delta(A)$ with the same barycentre
	$b:=\int p\,d\zeta^+(p)=\int p\,d\zeta^-(p)$.  Since $b\in\Delta(A)$, the set $\{a:b(a)>0\}$ is nonempty.  Because $r$
	has full support,
	\[
	\theta:=\min\Bigl\{1,\min_{a:b(a)>0}\frac{r(a)}{b(a)}\Bigr\}>0.
	\]
	Choose $0<\lambda<\theta$ and set
	\[
	\tau:=\frac{r-\lambda b}{1-\lambda}.
	\]
	Then $1-\lambda>0$, $\sum_a\tau(a)=1$, and $\tau(a)\ge0$ for every $a$: if $b(a)>0$ this follows from
	$\lambda<r(a)/b(a)$, while if $b(a)=0$ the numerator is $r(a)>0$.  Hence $\tau\in\Delta(A)$.  Define
	\[
	\nu:=\lambda\zeta^++(1-\lambda)\delta_\tau,
	\qquad
	\nu':=\lambda\zeta^-+(1-\lambda)\delta_\tau.
	\]
	Both are finite-support probability measures on $\Delta(A)$ with barycentre $\lambda b+(1-\lambda)\tau=r$, so
	$\nu,\nu'\in\mathcal M_r$.  Then $\nu-\nu'=\lambda(\zeta^+-\zeta^-)=(\lambda/c)\eta$, hence
	$\eta=(c/\lambda)(\nu-\nu')\in V_r$.  Thus $W\subset V_r$ and $V_r=W$.
Since $W$ depends only on the full-support property (not on the specific prior), $V_q=V_r=V_s=W$ for all full-support priors.

\emph{Case 1: Degenerate branch posterior ($|B|=1$).}  If $r=\delta_a$ for some action $a$, then
	$\mathcal M_{\delta_a}=\{\delta_{\delta_a}\}$ is a singleton: the only Bayes-plausible law with barycentre $\delta_a$
	is the point mass at $\delta_a$ itself.  Hence $V_r=\{0\}$, and every continuation experiment at prior $\delta_a$ gives
	the same posterior law $\delta_{\delta_a}$.  We set $F_{\delta_a}(\delta_{\delta_a})=0$.  For such branches, the term
	$m(o)\beta(q,r_o)F_{r_o}(\mu_{r_o,Q^o})$ is zero under any convention for $\beta(q,r_o)$.

For the aggregation formula below, assign any positive value to $\beta(q,\delta_a)$; the choice is not behaviourally
identified because the corresponding continuation value is always zero.


\emph{Case 2: Non-full-support branch posterior ($1<|B|<|A|$).}  If $r$ has support $B\subsetneq A$ with $|B|\ge 2$,
apply Lemma~\ref{lem:supprestrict}.  Continuation comparisons at prior $r$ reduce to those at $r|_B\in\Delta(B)$,
with representative $F_r:=F_{r|_B}^B\circ\pi_B$.

For any finite action set $C$, write
\[
W_C:=\{\sigma\in\operatorname{span}\Delta_f(\Delta(C)):\sigma(\Delta(C))=0,\ \textstyle\int_{\Delta(C)}p\,d\sigma(p)=0\}.
\]
Let $\jmath_B:\Delta(B)\to\Delta(A)$ be the face inclusion, extending each posterior by zero outside $B$, and let
$i_B:=(\jmath_B)_\#$ act on finite signed posterior laws.  By Lemma~\ref{lem:supprestrict}(iv),
$F_r^A=F_{r|_B}^B\circ\pi_B$, so on face directions
\[
L_r^A(i_B\sigma)=L_{r|_B}^B(\sigma)\qquad(\sigma\in W_B).
\]
Since $|B|\ge2$ and $r|_B$ is full support on $B$, local non-triviality and Lemma~\ref{lem:convnorm} give
\[
L_{r|_B}^B(\chi_{r|_B}-\delta_{r|_B})>0,
\]
so $L_{r|_B}^B\ne0$.

\emph{Full-to-face scalar.}  Fix any full-support $q\in\Delta(A)$.  The boundary posterior $r$ is reachable from $q$:
choose
\[
0<\varepsilon<\min\Bigl\{1,\min_{b\in B}\frac{q(b)}{r(b)}\Bigr\}
\]
and set $P_1(1\mid a)=\varepsilon r(a)/q(a)$, with $r(a)=0$ for $a\notin B$.  Then branch $1$ has positive probability
$\varepsilon$ and posterior $r$.  For any nonzero $\sigma\in W_B$, the tangent-space identification on the face $B$ gives
$\sigma=t(\nu-\nu')$ for some $t>0$ and $\nu,\nu'\in\mathcal M_{r|_B}^B$.  Implement $\nu,\nu'$ by continuation channels
on $B$ and extend them arbitrarily to $A$.  Lemma~\ref{lem:supprestrict} identifies their comparison at the boundary prior
$r$ with their comparison at $r|_B$.

Axiom~\textup{(A7)} applied to the branch reaching $r$, with all other branches fixed, gives the positive sign direction
for feasible differences.  The negative direction follows by swapping the two continuation laws, and equality follows by
applying the weak part of Axiom~\textup{(A7)} in both directions.  Extending by positive scalar multiples from feasible
differences to all of $W_B$, $L_q^A\circ i_B$ and $L_{r|_B}^B$ have the same sign partition on $W_B$.  Since
$L_{r|_B}^B\ne0$, both are nonzero and there is a unique $\beta_A(q,r)>0$ such that
\[
L_q^A(i_B\sigma)=\beta_A(q,r)L_{r|_B}^B(\sigma)\qquad(\sigma\in W_B).
\]
This is the ambient full-to-face coefficient used for boundary branches.

\emph{Case 3: Full-support branch posterior ($B=A$).}  This is the main case.  Suppose $r$ has full support on $A$.

\emph{Claim 1: aggregate change from branch variation.}  If a first-stage branch reaches posterior $r$ with probability
$m$, and the branch continuation law changes from $\nu$ to $\nu'$ (both in $\mathcal M_r$), then the aggregate posterior
law changes by $m(\nu-\nu')\in m\cdot V_r$.  This difference depends only on $(m,\nu,\nu',r)$---not on which experiment
produced the branch, nor on the other branches.

\emph{Proof of Claim 1.}  The aggregate changes from $m\nu+\rho$ to $m\nu'+\rho$, where $\rho$ collects the other
branches.  The difference is $m(\nu-\nu')$.

\emph{Claim 2: $L_q|_{V_r}$ is a positive scalar multiple of $L_r|_{V_r}$.}

\emph{Proof of Claim 2.}  We first show that every direction in $V_r$ is proportional to a feasible difference.
By the Tangent-space identification above, any nonzero $\sigma\in V_r$ can be written as
$\sigma=t(\nu-\nu')$ for some $t>0$ and $\nu,\nu'\in\mathcal M_r$.  Hence every nonzero $\sigma\in V_r$ is a positive
scalar multiple of some feasible difference.

By Axiom~\textup{(A7)} (branchwise monotonicity), if $\nu\succ_r\nu'$ (i.e.\ $F_r(\nu)>F_r(\nu')$, equivalently
$L_r(\nu-\nu')>0$), then the aggregate comparison ranks the $\nu$-continuation strictly higher.  Since the aggregate
value difference is $L_q(m(\nu-\nu'))=mL_q(\nu-\nu')$ with $m>0$, we have $L_q(\nu-\nu')>0$ whenever $L_r(\nu-\nu')>0$.

\emph{Reverse and equality directions.}  If $L_r(\nu-\nu')<0$, completeness of $\succeq_r$ gives $\nu'\succ_r\nu$, and
A7 applied to the reversed comparison gives $L_q(\nu'-\nu)>0$, i.e.\ $L_q(\nu-\nu')<0$.  If $L_r(\nu-\nu')=0$, then
$\nu\sim_r\nu'$, so both $\nu\succeq_r\nu'$ and $\nu'\succeq_r\nu$; applying A7 weakly in both directions gives
$L_q(\nu-\nu')\ge 0$ and $L_q(\nu-\nu')\le 0$, hence $L_q(\nu-\nu')=0$.

Since every nonzero $\sigma\in V_r$ is proportional to some feasible difference, the sign agreement extends from
feasible differences to all of $V_r$: for any $\sigma\in V_r$, $L_q(\sigma)>0$ iff $L_r(\sigma)>0$, $L_q(\sigma)<0$
iff $L_r(\sigma)<0$, and $L_q(\sigma)=0$ iff $L_r(\sigma)=0$.

Two nonzero linear functionals on a vector space that have the same positive half-space must be positive scalar
multiples of each other.  (If $L_r\ne 0$, then $\ker L_q=\ker L_r$ and they lie in the same half-space, so $L_q=cL_r$
for some $c>0$.)  Since $r$ has full support and local non-triviality gives $F_r(\chi_r)>F_r(\delta_r)=0$, we have
$L_r\ne 0$.  Hence there exists $\beta(q,r)>0$ such that
\[
L_q(\sigma)=\beta(q,r)\,L_r(\sigma)
\qquad\text{for all }\sigma\in V_r.
\]

\emph{Path-independence.}  Consider any first-stage experiment reaching posterior $r$ in some branch with probability
$m$.  Changing the branch continuation from $\nu$ to $\nu'$ changes the aggregate value by
\[
L_q(m(\nu-\nu'))=m\,\beta(q,r)\,L_r(\nu-\nu')=m\,\beta(q,r)\,(F_r(\nu)-F_r(\nu')).
\]
The coefficient $m\,\beta(q,r)$ depends only on $(q,r,m)$.  In particular, $\beta(q,r)$ is independent of which
experiment reached $r$ and independent of the other branches.

Comparing with the earlier notation $g(\nu)=\alpha F_r(\nu)+\gamma$: the value difference is
$g(\nu)-g(\nu')=\alpha(F_r(\nu)-F_r(\nu'))$, so $\alpha=m\,\beta(q,r)$ and hence $\alpha/m=\beta(q,r)$ universally.
The coefficient $\beta(q,r_{\bar o})$ is defined as the scalar from $L_q|_{V_{r_{\bar o}}}=\beta(q,r_{\bar o})L_{r_{\bar o}}|_{V_{r_{\bar o}}}$;
it depends only on $q$ and $r_{\bar o}$, not on the experiment that produced the branch.

\emph{Aggregation formula.}  Let
\[
\mu^0:=\sum_{o:m(o)>0}m(o)\delta_{r_o}=\mu_{q,P_1}.
\]
For each positive-probability branch let $B_o:=\operatorname{supp}(r_o)$.  If $B_o=A$, set
$v_o:=\mu_{r_o,Q^o}-\delta_{r_o}\in W_A$.  If $1<|B_o|<|A|$, set
\[
\bar v_o:=\pi_{B_o}(\mu_{r_o,Q^o})-\delta_{r_o|_{B_o}}\in W_{B_o},
\qquad
v_o:=i_{B_o}\bar v_o\in W_A.
\]
If $|B_o|=1$, set $v_o:=0$.  Then
\[
\mu_{q,P_1\triangleright\{Q^o\}}=\mu^0+\sum_{o:m(o)>0}m(o)v_o .
\]
Since $V_q=W_A$, affinity of $F_q$ gives
\[
F_q(\mu_{q,P_1\triangleright\{Q^o\}})-F_q(\mu^0)
=\sum_{o:m(o)>0}m(o)L_q^A(v_o).
\]
For a full-support branch $B_o=A$, Case~3 gives
\[
L_q^A(v_o)=\beta(q,r_o)L_{r_o}^A(v_o)
=\beta(q,r_o)F_{r_o}(\mu_{r_o,Q^o}).
\]
For a nondegenerate boundary branch $1<|B_o|<|A|$, Case~2 gives
\[
L_q^A(v_o)=L_q^A(i_{B_o}\bar v_o)
=\beta_A(q,r_o)L_{r_o|_{B_o}}^{B_o}(\bar v_o)
=\beta_A(q,r_o)F_{r_o}(\mu_{r_o,Q^o}).
\]
The last equalities use zero normalisation: Lemma~\ref{lem:convnorm} in full support, Lemma~\ref{lem:supprestrict} plus
Lemma~\ref{lem:convnorm} on the support face in the boundary case, and the singleton convention when $|B_o|=1$.  Therefore
\[
F_q(\mu_{q,P_1\triangleright\{Q^o\}})
=
F_q(\mu_{q,P_1})+\sum_{o:m(o)>0}m(o)\,\beta(q,r_o)\,F_{r_o}(\mu_{r_o,Q^o}),
\]
where $\beta(q,r_o)$ denotes the Case~3 coefficient when $B_o=A$, the Case~2 full-to-face coefficient when
$1<|B_o|<|A|$, and any positive singleton convention when $|B_o|=1$.
\end{proof}

\begin{lemma}[Branch-coefficient cocycle and normalised chain rule]\label{lem:chain}
Let the full-support and full-to-face coefficients be those constructed in Lemma~\ref{lem:branchagg}, applied in the
relevant ambient simplex or intrinsically on the relevant support face.  If $q$ has full support and $r,s$ are
nondegenerate priors with
\[
\operatorname{supp}(s)\subseteq\operatorname{supp}(r)\subseteq\operatorname{supp}(q),
\]
then the coefficients satisfy the cocycle identity
\[
\beta(q,s)=\beta(q,r)\beta(r,s).
\]
Hence $\beta(q,r)=a_q/a_r$ for positive scales $a_q$, with singleton scales fixed by convention.  The rescaled values
\[
\widehat F_q(\mu):=\frac{F_q(\mu)}{a_q}
\]
satisfy the exact chain rule: for every first-stage channel $P_1$ at full-support $q$ and every continuation profile
$\{Q^o\}_{o:m(o)>0}$,
\[
\widehat F_q(\mu_{q,P_1\triangleright\{Q^o\}})
=
\widehat F_q(\mu_{q,P_1})
\;+\;
\sum_{o:m(o)>0}m(o)\,\widehat F_{r_o}(\mu_{r_o,Q^o}).
\]
\end{lemma}

\begin{proof}

\emph{Full-support cocycle.}  First consider reachable full-support $q,r,s\in\operatorname{int}\Delta(A)$ with $|A|\ge2$.
For such priors, all three tangent spaces coincide.
By the tangent-space identification in Lemma~\ref{lem:branchagg}, $V_q=V_r=V_s=W$ for all full-support priors.
The full-support coefficient identities constructed there share a common domain $W$:
\[
L_q|_W=\beta(q,r)L_r|_W,
\qquad
L_r|_W=\beta(r,s)L_s|_W.
\]
Substituting the second into the first:
\[
L_q|_W=\beta(q,r)\beta(r,s)L_s|_W.
\]
Comparing with $L_q|_W=\beta(q,s)L_s|_W$ and using $L_s\ne 0$:
\[
\beta(q,s)=\beta(q,r)\beta(r,s).
\]

\emph{Reachability.}  The above uses the full-support coefficient construction for the pairs $(q,r)$, $(r,s)$, and
$(q,s)$.  Each invocation requires the pair to be reachable: the second prior must be a possible posterior of some
experiment at the first prior.
For any full-support $q,r$, construct a two-outcome experiment $P_1$ with $P_1(1\mid a)=\varepsilon r(a)/q(a)$ for
$0<\varepsilon<\min\{1,\min_a q(a)/r(a)\}$.  Outcome~1 has probability $\varepsilon$ and posterior~$r$; outcome~0 has posterior
$(q-\varepsilon r)/(1-\varepsilon)$.  Hence any full-support $r$ is reachable from any full-support $q$, so the construction
applies to all full-support pairs.  The cocycle therefore holds for all full-support triples.

Fix a full-support basepoint $q_0$, set $a_{q_0}=1$, and for full-support $q$ define
\[
a_q:=\frac{1}{\beta(q_0,q)}.
\]
By the full-support cocycle, this is equivalently $a_q=\beta(q,q_0)$, and
\[
\beta(q,r)=\frac{\beta(q,q_0)}{\beta(r,q_0)}=\frac{a_q}{a_r}.
\]

\emph{Boundary coefficients and face scales.}  If $r$ is nondegenerate with support $B\subsetneq A$, Lemma~\ref{lem:branchagg}
constructs the ambient full-to-face coefficient $\beta_A(q,r)>0$ by
\[
L_q^A i_B=\beta_A(q,r)L_{r|_B}^B .
\]
When the ambient action set is clear, write this coefficient as $\beta(q,r)$.

When nondegenerate boundary priors $r,s$ have common support $B$, $\beta_A(r,s)$ means the unique $c>0$ satisfying
\[
L_r^A(i_B\sigma)=c\,L_s^A(i_B\sigma)\qquad(\sigma\in W_B).
\]
This is a scalar on the common face, not on all of $W_A$.  Since
\[
L_r^A\circ i_B=L_{r|_B}^B,\qquad L_s^A\circ i_B=L_{s|_B}^B,
\]
and the intrinsic face scalar satisfies
$L_{r|_B}^B(\sigma)=\beta_B(r|_B,s|_B)L_{s|_B}^B(\sigma)$, we have
\[
\beta_A(r,s)=\beta_B(r|_B,s|_B).
\]

Fix the same full-support basepoint $q_0$ on $\Delta(A)$.  If $r_0$ is any nondegenerate face basepoint with support $B$,
then for every nondegenerate $r$ with support $B$,
\[
\beta_A(q_0,r)=\beta_A(q_0,r_0)\,\beta_B(r_0|_B,r|_B).
\]
Indeed, both sides are the coefficient multiplying $L_{r|_B}^B$ to obtain $L_{q_0}^A\circ i_B$, and
$L_{r|_B}^B\ne0$.  Define the ambient boundary scale by
\[
a_r:=\frac{1}{\beta_A(q_0,r)}.
\]
If the face scale is based at $r_0$ with $a_{r_0}:=1/\beta_A(q_0,r_0)$, then
\[
a_{r|_B}^B:=\frac{a_{r_0}}{\beta_B(r_0|_B,r|_B)}
=\frac{1}{\beta_A(q_0,r_0)\beta_B(r_0|_B,r|_B)}
=\frac{1}{\beta_A(q_0,r)}
=a_r.
\]
Thus the ambient and intrinsic scales agree: treating $r$ as a boundary posterior in $\Delta(A)$ or as a full-support prior
on $\Delta(B)$ yields the same normalised representative $\widehat F_r$.

\emph{Extension to boundary posteriors.}  If $r$ is nondegenerate with support $B\subsetneq A$, Lemma~\ref{lem:branchagg}
defines
$\beta(q,r)=\beta_A(q,r)$ by
\[
L_q^A i_B=\beta_A(q,r)L_{r|_B}^B .
\]
Using the full-support identity $L_q^A=\beta(q,q_0)L_{q_0}^A$ and restricting to $i_B(W_B)$,
\[
L_q^A i_B=\beta(q,q_0)L_{q_0}^A i_B
=\beta(q,q_0)\beta_A(q_0,r)L_{r|_B}^B.
\]
Comparing with the definition of $\beta_A(q,r)$ and using $L_{r|_B}^B\ne0$ gives
\[
\beta(q,r)=\beta(q,q_0)\beta(q_0,r)=\frac{a_q}{a_r},
\]
where $a_r=1/\beta(q_0,r)$ by the boundary-scale definition above.

\emph{Singleton scale conventions.}  If $r=\delta_a$ is a singleton prior, set $a_{\delta_a}:=a_{q_0}=1$ and
$\widehat F_{\delta_a}\equiv0$.  Define $\beta(q,\delta_a):=a_q/a_{\delta_a}$ as a convention.  No behavioural branch
coefficient is identified at a singleton target, but singleton terms in the aggregation and chain formulas are always zero.

\emph{Mixed-support cocycle.}  Let $r,s$ be nondegenerate with nested supports
$C:=\operatorname{supp}(s)\subseteq B:=\operatorname{supp}(r)$, and write $s|_B$ for $s$ viewed as a boundary prior on
the action set $B$.  Interpret
\[
\beta(r,s):=\beta_B(r|_B,s|_B),
\]
the intrinsic full-to-face coefficient on action set $B$.  Let $i_C^A:W_C\to W_A$ and $i_C^B:W_C\to W_B$ be the face
inclusions.  For every $\sigma\in W_C$,
\[
L_q^A i_C^A\sigma
=\beta(q,r)L_{r|_B}^B i_C^B\sigma
=\beta(q,r)\beta(r,s)L_{s|_C}^C\sigma.
\]
The direct full-to-face definition from Lemma~\ref{lem:branchagg} for the pair $(q,s)$ gives
\[
L_q^A i_C^A\sigma=\beta(q,s)L_{s|_C}^C\sigma.
\]
Since $L_{s|_C}^C\ne0$, $\beta(q,s)=\beta(q,r)\beta(r,s)$.  If
$\operatorname{supp}(s)\not\subseteq\operatorname{supp}(r)$, then $\beta(r,s)$ is not behaviourally defined; singleton
targets are also excluded from the behavioural cocycle because their scales are conventions only.

\emph{Chain rule.}  Applying Lemma~\ref{lem:branchagg} and substituting $\beta(q,r)=a_q/a_r$ gives
\[
F_q(\mu_{q,P_1\triangleright\{Q^o\}})
=
F_q(\mu_{q,P_1})+\sum_{o:m(o)>0}m(o)\,\frac{a_q}{a_{r_o}}\,F_{r_o}(\mu_{r_o,Q^o}).
\]
Dividing by $a_q$ and writing $\widehat F_q(\mu):=F_q(\mu)/a_q$ yields
\[
\widehat F_q(\mu_{q,P_1\triangleright\{Q^o\}})
=
\widehat F_q(\mu_{q,P_1})+\sum_{o:m(o)>0}m(o)\,\widehat F_{r_o}(\mu_{r_o,Q^o}).
\qedhere
\]
\end{proof}

\begin{lemma}[Coherent relabelling and face scales]\label{lem:facescales}
The scales produced by Lemma~\ref{lem:chain} can be chosen simultaneously across all finite action sets so that:
\begin{enumerate}
\item for every bijection $\pi:A\to A'$,
\[
a^{A'}_{q\pi}=a^A_q;
\]
\item if $\iota:B\hookrightarrow A$ is an injection with $|B|\ge2$, $r$ has full support on $B$, and $q$ has full support
on $A$, then the full-to-face branch coefficient for reaching the face $\iota(B)$ is
\begin{equation}
\beta_\iota(q,r)=\frac{a^A_q}{a^B_r}.
\tag{FS}\label{eq:facescale}
\end{equation}
\end{enumerate}
Singleton target scales remain a convention because every continuation value on a singleton fibre is zero.
\end{lemma}

\begin{proof}
For every nondegenerate finite action set $A$, first choose the full-support scales supplied by the cocycle part of
Lemma~\ref{lem:chain}, so that
\[
\beta_A(q,q')=\frac{a^A_q}{a^A_{q'}}
\qquad(q,q'\in\operatorname{int}\Delta(A)).
\]
They are unique up to multiplication by one positive constant depending on $A$.  Normalise initially by
$a^A_{u_A}=1$, where $u_A$ is uniform.  Exact relabelling invariance of the representatives implies relabelling
invariance of the uniquely determined branch coefficients.  Since bijections carry uniform priors to uniform priors,
the initial scales are exactly invariant under bijections.

Let $\iota:B\hookrightarrow A$ be an injection between nondegenerate finite sets.  By relabelling $B$ onto its image
$\iota(B)$ and using Corollary~\ref{cor:permutationinvariance}, it is enough to define the coefficient for the inclusion
of the face $\iota(B)$; the notation below keeps the original symbol $\iota$.  Let $i_\iota$ denote the induced map on
signed posterior-law tangent spaces.  For full-support $q\in\Delta(A)$ and $r\in\Delta(B)$, let
$\beta_\iota(q,r)>0$ be the full-to-face scalar characterised by
\[
L_q^A\circ i_\iota=\beta_\iota(q,r)L_r^B
\quad\hbox{on }W_B.
\]
The number
\begin{equation}
K(\iota):=\beta_\iota(q,r)\frac{a^B_r}{a^A_q}
\tag{K}\label{eq:embeddingdefect}
\end{equation}
is independent of both $q$ and $r$.  Independence of $q$ follows by comparing two full-support ambient priors and using
$L_q^A=(a_q^A/a_{q'}^A)L_{q'}^A$ on $W_A$.  Independence of $r$ follows from same-face compatibility: if $r,s$ have
full support on $B$, then $L_r^B=(a_r^B/a_s^B)L_s^B$, and hence
$\beta_\iota(q,s)=\beta_\iota(q,r)a_r^B/a_s^B$.

The defects multiply under composition.  If
$C\xhookrightarrow{\jmath}B\xhookrightarrow{\iota}A$ and all three sets in the comparison are nondegenerate, then
restricting first to $B$ and then to $C$ gives, for full-support $q,r,s$ on $A,B,C$ respectively,
\[
\beta_{\iota\circ\jmath}(q,s)=\beta_\iota(q,r)\beta_\jmath(r,s).
\]
Substituting the definition of $K$ and cancelling the intermediate scale $a_r^B$ gives
\begin{equation}
K(\iota\circ\jmath)=K(\iota)K(\jmath).
\tag{KC}\label{eq:embeddingcocycle}
\end{equation}
Exact relabelling invariance shows that $K(\iota)$ depends only on $n:=|A|$ and $m:=|B|$; write it as $K_{n,m}$.
Bijections give $K_{n,n}=1$, and
\[
K_{n,\ell}=K_{n,m}K_{m,\ell}
\qquad(2\le\ell\le m\le n).
\]
Put $t_n:=K_{n,2}$, with $t_2=1$.  Then $K_{n,m}=t_n/t_m$.  Finally replace every scale on an $n$-point action set by
\[
\widetilde a_q^A:=t_n a_q^A.
\]
Under this change, the defect in~\eqref{eq:embeddingdefect} becomes
\[
\widetilde K_{n,m}=K_{n,m}\frac{t_m}{t_n}=1.
\]
The rescaling is common to all priors on a fixed action set, so it leaves all within-set cocycle ratios unchanged; because
$t_n$ depends only on cardinality, it also preserves exact relabelling invariance.  Dropping tildes gives
\eqref{eq:facescale}.  Singleton target scales may be set arbitrarily because all singleton continuation terms are zero.
\end{proof}

\begin{lemma}[Interaction collapse and universal chain scale]\label{lem:scalecoherence}
The universal interaction coefficient in~\eqref{eq:QA} is zero.  Moreover, there exists $a>0$ such that the coherently
chosen chain scales satisfy $a_q=a$ for every nondegenerate full-support prior on every finite action set.  The same
value may be assigned to singleton priors.  Consequently,
\begin{equation}
F_{p\otimes r}(\mu\otimes\nu)=F_p(\mu)+F_r(\nu)
\tag{PA}\label{eq:PArecovered}
\end{equation}
and the branch formula has a universal chain scale.
\end{lemma}

\begin{proof}
Define, for every finite prior $q$ read on its support,
\[
H(q):=F_q(\chi_q),
\qquad
Z(q):=1+\kappa H(q).
\]
For a nondegenerate prior, local non-triviality gives $H(q)>0$, and~\eqref{eq:QAslope} gives
\begin{equation}
Z(q)>0.
\tag{Z+}\label{eq:Zpositive}
\end{equation}
For a singleton, $H(q)=0$ and $Z(q)=1$.

\emph{Step 1: product revelation links the chain scale to $Z$.}
Let $q\in\Delta(A)$ and $r\in\Delta(B)$ be nondegenerate and full support.  Quasi-additivity applied to full revelation
gives
\begin{equation}
H(q\otimes r)=H(q)+H(r)+\kappa H(q)H(r).
\tag{QH}\label{eq:QH}
\end{equation}
Compute the same value sequentially.  First reveal the $A$-coordinate and then reveal the $B$-coordinate in every branch.
The first-stage value is $F_{q\otimes r}(\chi_q\otimes\delta_r)=H(q)$ by~\eqref{eq:QA}.  After observing $a$, the
posterior is $e_a\otimes r$ on the face $\{a\}\times B$; support restriction and relabelling identify its continuation
full-revelation value with $H(r)$.  For the injection
\[
\iota_a:B\hookrightarrow A\times B,\qquad b\mapsto(a,b),
\]
Lemma~\ref{lem:facescales} gives the branch coefficient
\[
\beta_{\iota_a}(q\otimes r,r)=\frac{a^{A\times B}_{q\otimes r}}{a^B_r}.
\]
The branch weights are $q(a)$ and the same coefficient and continuation value appear in every branch, so
$\sum_a q(a)=1$.  Lemma~\ref{lem:branchagg} therefore yields
\begin{equation}
H(q\otimes r)=H(q)+\frac{a_{q\otimes r}}{a_r}H(r).
\tag{SA}\label{eq:seqA}
\end{equation}
Comparing~\eqref{eq:QH} and~\eqref{eq:seqA}, and using $H(r)>0$,
\begin{equation}
\frac{a_{q\otimes r}}{a_r}=Z(q).
\tag{R1}\label{eq:ratio1}
\end{equation}
Revealing $B$ first gives symmetrically
\begin{equation}
\frac{a_{q\otimes r}}{a_q}=Z(r).
\tag{R2}\label{eq:ratio2}
\end{equation}
Hence $a_rZ(q)=a_qZ(r)$, so there is a universal constant $C>0$ such that
\begin{equation}
a_q=CZ(q)
\tag{AS}\label{eq:ascaleZ}
\end{equation}
for every nondegenerate full-support prior on every finite action set.  Assign $a_q:=C$ at singleton priors, where
$Z(q)=1$ and all continuation values are zero.

If $\kappa=0$, then $Z\equiv1$, equation~\eqref{eq:ascaleZ} already gives universal scales, and~\eqref{eq:QA} is exact
product additivity.  For the remaining argument suppose $\kappa\ne0$.

\emph{Step 2: the induced weight equation.}
Let $q$ be full support on $A$, let $\mathcal P=\{A_1,\ldots,A_m\}$ be a partition into nonempty cells, and put
\[
p_i:=q(A_i),
\qquad
q^i:=q(\cdot\mid A_i),
\qquad
p:=(p_1,\ldots,p_m).
\]
Let $G_{\mathcal P}$ reveal the block.  Undetectable-distinctions neutrality gives
$(q,G_{\mathcal P})\sim(p,\Id_{\{1,\ldots,m\}})$.  The early cross-prior representation
Lemma~\ref{lem:blockbridge}, importantly for the \emph{unscaled} $F$'s, therefore gives
\begin{equation}
F_q(\mu_{q,G_{\mathcal P}})=H(p).
\tag{BG}\label{eq:blockvalue}
\end{equation}
Append full revelation in every branch.  The compound posterior law is $\chi_q$, so the chain formula and
\eqref{eq:ascaleZ} give
\begin{equation}
H(q)=H(p)+Z(q)\sum_{i=1}^m p_i\frac{H(q^i)}{Z(q^i)}.
\tag{GR}\label{eq:groupingH}
\end{equation}
Define
\[
w(q):=\frac1{Z(q)}>0.
\]
Using $\kappa H(x)=Z(x)-1$ in~\eqref{eq:groupingH} gives
\[
\frac{Z(p)}{Z(q)}=\sum_{i=1}^m p_i\frac1{Z(q^i)}.
\]
Equivalently,
\begin{equation}
w(q)=w(p)\sum_{i=1}^m p_iw(q^i).
\tag{W}\label{eq:weight}
\end{equation}
No continuity assertion is needed here.

\emph{Step 3: a two-grouping argument forces $w\equiv1$.}
All distributions in this paragraph are read on their positive-probability supports, so the weight equation applies after
support restriction.  First,~\eqref{eq:weight} applied to a product distribution gives
\begin{equation}
w(u\otimes v)=w(u)w(v),
\tag{M}\label{eq:wmult}
\end{equation}
because the first-coordinate block probabilities are $u$ and every conditional distribution is $v$.

Take arbitrary finite probability vectors $u$ and $v$, possibly on different alphabets, and write
\[
x:=w(u),\qquad y:=w(v),\qquad p:=(1/2,1/2).
\]
On the labelled disjoint union $(U\times U)^0\sqcup(V\times V)^1$, form the distribution
\[
T:=\tfrac12(u\otimes u)^0\;\sqcup\;\tfrac12(v\otimes v)^1.
\]
Grouping first by the top block and using~\eqref{eq:wmult},
\begin{equation}
w(T)=w(p)\frac{x^2+y^2}{2}.
\tag{E1}\label{eq:twoeval1}
\end{equation}
Now group first by the pair consisting of the top-block label and the first within-block coordinate.  Its marginal is
\[
S:=\tfrac12u^0\;\sqcup\;\tfrac12v^1,
\]
and~\eqref{eq:weight} gives
\[
w(S)=w(p)\frac{x+y}{2}.
\]
Conditional on a point in the first part of $S$, the remaining coordinate has distribution $u$; conditional on a point in
the second part, it has distribution $v$.  A second application of~\eqref{eq:weight} therefore gives
\begin{equation}
w(T)=w(S)\frac{x+y}{2}
=w(p)\left(\frac{x+y}{2}\right)^2.
\tag{E2}\label{eq:twoeval2}
\end{equation}
Since $w(p)>0$, comparison of~\eqref{eq:twoeval1} and~\eqref{eq:twoeval2} yields
\[
\frac{x^2+y^2}{2}=\left(\frac{x+y}{2}\right)^2,
\]
so $(x-y)^2=0$.  As $u$ and $v$ were arbitrary, $w$ is constant on all finite probability vectors.  At a singleton prior,
$H=0$, hence $w=1$.  Therefore $w(q)\equiv1$.

Thus $Z(q)=1$ and $\kappa H(q)=0$ for every $q$.  Local non-triviality gives a nondegenerate $q$ with $H(q)>0$, hence
$\kappa=0$, contradicting the temporary alternative unless the interaction coefficient is zero.  Equation
\eqref{eq:ascaleZ} now gives $a_q=C$ universally.  Finally~\eqref{eq:QA} reduces to exact product additivity
\eqref{eq:PArecovered}.
\end{proof}

\begin{lemma}[Entropy reduction and Faddeev]\label{lem:faddeevsketch}
By Lemma~\ref{lem:scalecoherence}, $a_q\equiv a$ is universal.  Define
\[
\widehat F_q:=F_q/a,
\qquad
H(q):=\widehat F_q(\chi_q),
\]
the rescaled value of full revelation.  Then, for every channel $P$,
\[
\widehat F_q(\mu_{q,P})
=
H(q)-\sum_{o:m_{q,P}(o)>0}m_{q,P}(o)\,H(r_o^{\,q,P}).
\tag{ER}
\]
Moreover $H$ satisfies Faddeev's recursion, so
\[
H(q)=\alpha \Sh(q)
\]
for some $\alpha>0$.  Consequently,
\[
\widehat F_q(\mu_{q,P})=\alpha I_{q,P}(A;O).
\]
\end{lemma}

\begin{proof}
\emph{Entropy-reduction formula.}  Boundary values of $H$ are read through support restriction: if $q$ has support
$B$, then $H(q):=H(q|_B)$, and $H(\delta_a)=0$.  Let $P:A\to\Delta(O)$ be any channel at full-support prior $q$.  Append
full revelation $\Id_A:A\to\Delta(A)$ after $P$: the compound $P\triangleright\{\Id_A\}_o$ (using $\Id_A$ in every branch)
has $\mu_{r_o,\Id_A}=\chi_{r_o}$ and the overall posterior law is full revelation $\chi_q$.  Applying the chain rule
(Lemma~\ref{lem:chain}) gives
\[
\widehat F_q(\chi_q)
=
\widehat F_q(\mu_{q,P})+\sum_{o:m(o)>0}m(o)\,\widehat F_{r_o}(\chi_{r_o}).
\]
Hence
\[
\widehat F_q(\mu_{q,P})=H(q)-\sum_{o:m(o)>0}m(o)\,H(r_o).
\]
If $q$ is on the boundary, restrict $q$ and $P$ to $B=\operatorname{supp}(q)$, apply the full-support formula on
$\Delta(B)$, and pull the value and comparison back by Lemma~\ref{lem:supprestrict}; zero-probability rows and outcomes
do not affect either side of the displayed identity.

\emph{Scaled cross-prior bridge.}
Lemma~\ref{lem:blockbridge} was proved for the unscaled values $F_q$.  Since Lemma~\ref{lem:scalecoherence} has made
$a_q\equiv a$ universal, it is equivalently valid for
\[
V(q,P):=\widehat F_q(\mu_{q,P})=F_q(\mu_{q,P})/a:
\]
for all finite priors and channels,
\[
q^0\succeq_{P\sqcup Q}r^1
\quad\Longleftrightarrow\quad
V(q,P)\ge V(r,Q).
\]

\emph{Faddeev recursion.}  First let $q$ be full support and let $\mathcal P=\{A_1,\dots,A_m\}$ be a partition into
nonempty cells.  Then $p_i:=q(A_i)>0$ for every $i$.  Let $G_{\mathcal P}:A\to\Delta(\{1,\dots,m\})$ reveal only the
block index: $G_{\mathcal P}(i\mid a)=\1_{a\in A_i}$.  Write
\[
p_i:=q(A_i),\qquad q^i:=q(\cdot\mid A_i),
\]
and let $p=(p_1,\dots,p_m)\in\Delta(\{1,\dots,m\})$.  Consider the deterministic kernel
$S:A\to\Delta(\{1,\dots,m\})$ given by
$S(i\mid a)=\1_{a\in A_i}$, so $qS=p$.

By Lemma~\ref{lem:neutrality}, since $G_{\mathcal P}$ is $S$-measurable,
\[
(q,G_{\mathcal P})\sim(p,\Id_{\{1,\dots,m\}})
\]
in the two-block comparison.  Applying the scaled form of Lemma~\ref{lem:blockbridge} to the forward comparison gives
$V(q,G_{\mathcal P})\ge V(p,\Id_{\{1,\dots,m\}})$.  Applying the reverse-block form from Lemma~\ref{lem:blockcoh} and then
the same bridge to the reverse comparison gives $V(p,\Id_{\{1,\dots,m\}})\ge V(q,G_{\mathcal P})$.  Therefore
\[
\widehat F_q(\mu_{q,G_{\mathcal P}})
=\widehat F_p(\chi_p)=H(p).
\]
The entropy-reduction formula applied to $G_{\mathcal P}$ gives
\[
H(q)=\widehat F_q(\chi_q)
=\widehat F_q(\mu_{q,G_{\mathcal P}})+\sum_{i=1}^m p_i\widehat F_{q^i}(\chi_{q^i})
=H(p)+\sum_{i=1}^m p_iH(q^i).
\]
This is Faddeev's recursion.  For boundary $q$, restrict to $\operatorname{supp}(q)$, apply the same full-support recursion
to the nonempty positive-mass intersections, and pull back by Lemma~\ref{lem:supprestrict}; equivalently, the recursion is
read over positive-mass blocks only.

We now verify that $H$ is continuous on each $\Delta_n:=\{q\in\Delta(A):|A|=n\}$.

\emph{Binary continuity via erasure calibrators.}
Define the binary entropy function $h:[0,1]\to\R$ by $h(t):=H(t,1-t)$.  We show $h$ is continuous using
prior-independent calibrators.  By Lemma~\ref{lem:convnorm}, support restriction, and $a>0$, $H(q)\ge0$ for every $q$,
so $h(t)\ge0$ on $[0,1]$.  If $c<0$, then
\[
\{t:h(t)\ge c\}=[0,1],
\qquad
\{t:h(t)\le c\}=\varnothing,
\]
both closed.  Now fix $c\ge0$.  By local non-triviality, there exists a binary prior $\theta=(s,1-s)$ with
$u:=H(\theta)>0$.  Choose $n\in\mathbb{N}$ and $\lambda\in[0,1]$ such that $c=\lambda nu$
(take $n$ large enough that $c\le nu$, then set $\lambda:=c/(nu)$).

Let $w_n:=\theta^{\otimes n}$ be the $n$-fold product prior on $B_n:=\{0,1\}^n$.  By product additivity,
$H(w_n)=nH(\theta)=nu$.  Define the erasure experiment $E_{n,\lambda}:B_n\to\Delta(\{*,1,\dots,2^n\})$ which
reveals the state fully with probability $\lambda$ and reveals nothing with probability $1-\lambda$.  By the
entropy-reduction formula,
\[
\widehat F_{w_n}(\mu_{w_n,E_{n,\lambda}})=\lambda H(w_n)=\lambda nu=c.
\]
Now consider the block action space $A_c:=\{0,1\}\sqcup B_n$ and the block-diagonal channel
$K_c:=\Id_{\{0,1\}}\sqcup E_{n,\lambda}$.  For $t\in[0,1]$, let $q_t^0$ be the prior $(t,1-t)$ supported on
block $\{0,1\}$, and let $w_n^1$ be $w_n$ supported on block $B_n$.  By the scaled form of
Lemma~\ref{lem:blockbridge} (which extends to boundary priors via Lemma~\ref{lem:supprestrict}),
\[
q_t^0\succeq_{K_c}w_n^1\quad\Longleftrightarrow\quad V(q_t,\Id_{\{0,1\}})=h(t)\ge c=V(w_n,E_{n,\lambda}).
\]
Since $K_c$ is fixed and the map $t\mapsto q_t^0$ is continuous in $\ell_1$, Axiom~\textup{(A2)} implies that the
set $\{t\in[0,1]:h(t)\ge c\}$ is closed.  Similarly, $\{t\in[0,1]:h(t)\le c\}$ is closed.  Thus all real superlevel and
sublevel sets of $h$ are closed, so $h$ is both upper and lower semicontinuous on $[0,1]$, hence continuous.

\emph{Continuity on higher simplices by induction.}
For $n\ge3$, write any $q\in\Delta_n$ as $q=(q_1,(1-q_1)\bar q)$ where
$\bar q:=(q_2/(1-q_1),\dots,q_n/(1-q_1))\in\Delta_{n-1}$ for $q_1<1$.  The Faddeev recursion with the
binary partition $\{\{1\},\{2,\dots,n\}\}$ gives
\[
H(q)=h(q_1)+(1-q_1)H(\bar q).
\]
On the region $0<q_1<1$, this is continuous: $h$ is continuous by binary continuity, $H|_{\Delta_{n-1}}$ is
continuous by the induction hypothesis, and $q\mapsto(q_1,\bar q)$ is continuous.  If $q_1=0$, support restriction,
equivalently the positive-mass-block convention in the recursion, gives $H(q)=H(\bar q)$ and $h(0)=0$,
so continuity at that face follows from the induction hypothesis.  At $q_1=1$, the induction hypothesis and compactness
of $\Delta_{n-1}$ imply $H$ is bounded on $\Delta_{n-1}$, so
$(1-q_1)H(\bar q)\to0$ as $q_1\to1$.  Thus $H(q)\to h(1)=0=H(\delta_1)$.  The other boundary faces are handled
by support restriction, relabelling on the support, and pulling back.  This completes the induction.

We have verified the hypotheses of the strong-additivity form of Faddeev's finite entropy characterisation, as stated
for example by Baez--Fritz--Leinster \cite[Theorem~6]{BaezFritzLeinster2011}, following Faddeev's original theorem
\cite{Faddeev1956}: a function on finite probability vectors that is continuous on each closed simplex, invariant under
bijections, expansive under adjoining zero-probability states, zero on point masses, and strongly additive is a
nonnegative multiple of Shannon entropy.  Here invariance under bijections follows from
Corollary~\ref{cor:permutationinvariance}, with boundary cases obtained by restricting to the support and pulling back;
expansibility follows from support restriction; and the grouping recursion above is precisely strong additivity, read
over positive-mass blocks.  Faddeev's theorem therefore gives $H(q)=\alpha\Sh(q)$ for some $\alpha\ge0$.
Local non-triviality (Axiom~\textup{(A1)}) forces $\alpha>0$.

\emph{Mutual information.}  Substituting $H=\alpha\Sh$ into the entropy-reduction formula yields
\[
\widehat F_q(\mu_{q,P})
=\alpha\Sh(q)-\sum_{o:m_{q,P}(o)>0}m_{q,P}(o)\alpha\Sh(r_o^{\,q,P})
=\alpha\bigl(\Sh(q)-\E[\Sh(r_O)]\bigr)
=\alpha I_{q,P}(A;O).
\qedhere
\]
\end{proof}

\begin{lemma}[Fixed-environment representation]\label{lem:globalsketch}
For every finite channel $P:A\to\Delta(O)$ and all $q,q'\in\Delta(A)$,
\[
q\succeq_P q'
\quad\Longleftrightarrow\quad
I_{q,P}(A;O)\ge I_{q',P}(A;O).
\]
\end{lemma}

\begin{proof}
Let
\[
V(s,R):=\widehat F_s(\mu_{s,R}),
\]
with boundary priors interpreted by Lemma~\ref{lem:supprestrict}.  Lemma~\ref{lem:blockbridge}, divided by the universal
scale from Lemma~\ref{lem:scalecoherence}, gives the cross-prior block-comparison bridge: for finite action sets,
priors $s,t$, and channels $R,S$,
\[
s^0\succeq_{R\sqcup S}t^1
\quad\Longleftrightarrow\quad
V(s,R)\ge V(t,S),
\]
including boundary priors by support restriction.  Lemma~\ref{lem:faddeevsketch} gives, for every finite
prior/channel pair,
\[
V(s,R)=\alpha I_{s,R}
\]
with $\alpha>0$.

Fix $P:A\to\Delta(O)$ and $q,q'\in\Delta(A)$.  By the first clause of Axiom~\textup{(A3)},
\[
q\succeq_P q'
\quad\Longleftrightarrow\quad
q^0\succeq_{P\sqcup P}(q')^1.
\]
Applying this bridge with $(s,R)=(q,P)$ and $(t,S)=(q',P)$ gives
\[
q^0\succeq_{P\sqcup P}(q')^1
\quad\Longleftrightarrow\quad
V(q,P)\ge V(q',P).
\]
Substituting $V=\alpha I$ and using $\alpha>0$ yields
\[
q\succeq_P q'
\quad\Longleftrightarrow\quad
I_{q,P}(A;O)\ge I_{q',P}(A;O).
\]

If $q$ or $q'$ is on the boundary, the bridge is read after restricting each pair to its positive-probability support
and then pulling the comparison back by Lemma~\ref{lem:supprestrict}.  Deleting
zero-probability actions does not change mutual information, and the Axiom~\textup{(A3)} equivalence remains the ambient
fixed-channel comparison for $P$.
\end{proof}

\section*{Worked example}

\begin{example}[Donor attribution]\label{ex:donor-attribution}
A donor can route a gift either to a named project $N$ or to a pooled fund $F$:
\[
A=\{N,F\}.
\]
The public record is binary.  It either contains a project-linked report $L$ or only an aggregate report $U$:
\[
O=\{L,U\}.
\]
Let
\[
P(L\mid N)=\alpha,\qquad P(L\mid F)=\beta,
\qquad 0\le \beta\le \alpha\le 1.
\]
When $\beta<\alpha$, the named project is more likely to leave a project-linked trace, but pooled funding may also
occasionally be linked and named projects may sometimes disappear into aggregate reporting.  Write $q(N)=p$ and
$q(F)=1-p$.  The probability of a linked report is
\[
m_{q,P}(L)=\beta+(\alpha-\beta)p.
\]
Therefore
\[
I_{q,P}(A;O)
=
h\!\big(\beta+(\alpha-\beta)p\big)
-p\,h(\alpha)-(1-p)h(\beta),
\]
where
\[
h(x):=-x\log x-(1-x)\log(1-x)
\]
is binary entropy, with the convention $0\log0=0$.

The formula isolates attribution rather than payoff.  If $\alpha=\beta$, the public record is independent of the donor's
act and traceability is zero for every $q$.  If $\alpha=1$ and $\beta=0$, the record fully reveals the realised act and
traceability is $h(p)$.  If $p=0$ or $p=1$, traceability is again zero, because there is no uncertainty about the act for
the report to resolve.  Thus traceability comes from statistical separation in the reporting distributions together with
uncertainty about which act was realised.
\end{example}

\section*{Discussion}

Three points are worth emphasising.

First, the theorem is \emph{within}-environment in its behavioural content.  The decision-maker chooses $q$ while the
channel $P$ indexes the decision problem.  Cross-problem-looking comparisons are routed through block comparison
environments, so they remain ordinary comparisons between lotteries in one fixed environment.  This is analogous to the
role played by statewise coherence principles in classic decision theory: they constrain how one and the same behavioural
criterion is represented across transformed problems without turning the environment into an object of choice.

Second, the product and branch axioms align the fibrewise cardinal scales before the entropy/Faddeev identification.
Axiom~\textup{(A8)} (independent-background separability) identifies the coordinate slice orders on product posterior
domains.  Separate affinity then yields the quasi-additive product form in Lemma~\ref{lem:coherentnorm}; branch
aggregation, face-scale coherence, and the two-grouping argument in Lemma~\ref{lem:scalecoherence} eliminate the possible
interaction term.

Third, the representation is not an expected-utility theorem.  If two mechanisms induce the same lottery on final payoffs
but different joint laws of action and outcome, expected utility is indifferent while the present model need not be.  In that
sense the model is non-consequentialist over payoff outcomes alone.
This fixed-environment formulation places the paper close in spirit to the behavioural literatures on preference
for flexibility \cite{Kreps1979}, temptation and self-control \cite{GulPesendorfer2001}, and the intrinsic value of
decision rights \cite{BartlingFehrHerz2014}, while the mathematical backbone remains the information-theoretic and
statistical-experiment traditions of Blackwell \cite{Blackwell1953} and Faddeev \cite{Faddeev1956}.
One can, however, combine the traceability term with ordinary
utility over final outcomes:
\[
U_P(q)=\E_{q,P}[u(O)]+\lambda I_{q,P}(A;O).
\]
The fixed-environment theorem identifies the second term.

\sloppy
\begin{thebibliography}{99}
\raggedright

\bibitem{BartlingFehrHerz2014}
Bj\"orn Bartling, Ernst Fehr, and Holger Herz.
\newblock The intrinsic value of decision rights.
\newblock {\em Econometrica}, 82(6):2005--2039, 2014.

\bibitem{BaezFritzLeinster2011}
John C. Baez, Tobias Fritz, and Tom Leinster.
\newblock A characterization of entropy in terms of information loss.
\newblock {\em Entropy}, 13(11):1945--1957, 2011.

\bibitem{Blackwell1953}
David Blackwell.
\newblock Equivalent comparisons of experiments.
\newblock {\em Annals of Mathematical Statistics}, 24(2):265--272, 1953.

\bibitem{Debreu1960}
G\'erard Debreu.
\newblock Topological methods in cardinal utility theory.
\newblock In Kenneth J. Arrow, Samuel Karlin, and Patrick Suppes, editors,
  {\em Mathematical Methods in the Social Sciences, 1959}, pages 16--26.
  Stanford University Press, 1960.

\bibitem{Faddeev1956}
D. K. Faddeev.
\newblock On the concept of entropy of a finite probabilistic scheme.
\newblock {\em Uspekhi Matematicheskikh Nauk}, 11(1):227--231, 1956.
\newblock In Russian.

\bibitem{GulPesendorfer2001}
Faruk Gul and Wolfgang Pesendorfer.
\newblock Temptation and self-control.
\newblock {\em Econometrica}, 69(6):1403--1435, 2001.

\bibitem{HersteinMilnor1953}
I. N. Herstein and John Milnor.
\newblock An axiomatic approach to measurable utility.
\newblock {\em Econometrica}, 21(2):291--297, 1953.

\bibitem{Kreps1979}
David M. Kreps.
\newblock A representation theorem for preference for flexibility.
\newblock {\em Econometrica}, 47(3):565--577, 1979.

\bibitem{MatejkaMcKay2015}
Filip Mat\v{e}jka and Alisdair McKay.
\newblock Rational inattention to discrete choices: A new foundation for the multinomial logit model.
\newblock {\em American Economic Review}, 105(1):272--298, 2015.

\bibitem{Sims2003}
Christopher A. Sims.
\newblock Implications of rational inattention.
\newblock {\em Journal of Monetary Economics}, 50(3):665--690, 2003.

\end{thebibliography}

\end{document}
